%=======================================================================
%  Projeto Final - Teoria da Computacao e Compiladores
%  DRESS-GL: artigo final, versao 3 (dois hospedeiros e ponte CoppeliaSim)
%  Template: IEEE Conference Template (Official) - compilar com pdfLaTeX
%=======================================================================
\documentclass[conference]{IEEEtran}
\IEEEoverridecommandlockouts

\usepackage{cite}
\usepackage{amsmath,amssymb,amsfonts}
\usepackage{graphicx}
\usepackage{textcomp}
\usepackage{xcolor}
\usepackage{array}
\usepackage{tabularx}
\usepackage[utf8]{inputenc}
\usepackage[T1]{fontenc}
\usepackage[brazilian]{babel}
\usepackage{microtype}
\usepackage{listings}
\usepackage{booktabs}
\usepackage{url}
\usepackage{adjustbox}
\usepackage{tikz}
\usetikzlibrary{automata,positioning,arrows.meta,babel}

% Aspas retas aparecem no codigo da DSL; desliga o atalho " do babel.
\AtBeginDocument{\shorthandoff{"}}

% Quebra de URL apenas em barras, pontos e hifens.
\def\UrlBreaks{\do\/\do.\do-}

% Evita linhas que invadem a margem.
\setlength{\emergencystretch}{1.5em}

% Mais espaco para floats (evita tabelas depois das referencias).
\renewcommand{\topfraction}{0.9}
\renewcommand{\bottomfraction}{0.8}
\renewcommand{\textfraction}{0.07}
\renewcommand{\floatpagefraction}{0.8}
\renewcommand{\dbltopfraction}{0.9}
\renewcommand{\dblfloatpagefraction}{0.8}
\setcounter{topnumber}{3}
\setcounter{dbltopnumber}{3}
\setcounter{totalnumber}{5}

% Coluna X alinhada a esquerda (tabularx) e cabecalho em duas linhas.
\newcolumntype{L}{>{\raggedright\arraybackslash}X}
\newcommand{\cab}[1]{\begin{tabular}[b]{@{}c@{}}#1\end{tabular}}
\renewcommand{\lstlistingname}{Listagem}
% Listagem + legenda em bloco indivisivel, com legenda na largura da coluna.
\newenvironment{blocolst}%
  {\par\medskip\noindent\begin{minipage}{\columnwidth}}%
  {\end{minipage}\par\medskip}
% \legendalst{rotulo}{texto}: numera pelo mesmo contador do listings.
\newcommand{\legendalst}[2]{%
  \par\vspace{2pt}%
  \refstepcounter{lstlisting}\label{#1}%
  {\footnotesize\lstlistingname~\thelstlisting.\enspace #2\par}}
\renewcommand{\IEEEkeywordsname}{Palavras-chave}

%----- realce de sintaxe -----
\lstdefinelanguage{DRESSGL}{
  morekeywords={POLICY,RULE,PRIORITY,GIVEN,WHEN,THEN,ASSERT,END,
                BEFORE,AFTER,AROUND,AND,OR,NOT,IN,SET,SCALE,INJECT,
                SUPPRESS,HIGHLIGHT,MARK,STALE,ARGMAX,ARGMIN,WITHIN,
                BY,view,pass,instances,resolution,budget,reserve,cooldown},
  morekeywords=[2]{ms,Hz,frames},
  morecomment=[l]{\#},
  morestring=[b]",
  sensitive=true
}
\lstset{
  basicstyle=\ttfamily\scriptsize,
  keywordstyle=\bfseries,
  keywordstyle=[2]\itshape,
  commentstyle=\itshape\color{gray},
  stringstyle=\color{black!70},
  numbers=left, numberstyle=\tiny, numbersep=4pt,
  frame=single, framesep=3pt,
  columns=fullflexible, keepspaces=true,
  breaklines=true, breakatwhitespace=true,
  postbreak=\mbox{\textcolor{gray}{$\hookrightarrow$}\space},
  captionpos=b, xleftmargin=12pt,
  linewidth=\dimexpr\columnwidth-12pt\relax,
  showstringspaces=false,
  upquote=true,
  extendedchars=true,
  literate={á}{{\'a}}1 {à}{{\`a}}1 {â}{{\^a}}1 {ã}{{\~a}}1
           {é}{{\'e}}1 {ê}{{\^e}}1 {í}{{\'\i}}1
           {ó}{{\'o}}1 {ô}{{\^o}}1 {õ}{{\~o}}1 {ú}{{\'u}}1 {ü}{{\"u}}1
           {ç}{{\c c}}1
           {Á}{{\'A}}1 {À}{{\`A}}1 {Â}{{\^A}}1 {Ã}{{\~A}}1
           {É}{{\'E}}1 {Ê}{{\^E}}1 {Í}{{\'I}}1
           {Ó}{{\'O}}1 {Ô}{{\^O}}1 {Õ}{{\~O}}1 {Ú}{{\'U}}1 {Ç}{{\c C}}1
}
\lstdefinestyle{saida}{numbers=none, xleftmargin=0pt, linewidth=\columnwidth,
  basicstyle=\ttfamily\tiny, keywordstyle=, stringstyle=,
  breakindent=1.5em}
\lstdefinestyle{dsl}{language=DRESSGL,
  basicstyle=\ttfamily\fontsize{6.5}{7.8}\selectfont}
  \DeclareMathSizes{6.5}{6.5}{5}{5}
\begin{document}

\title{DRESS-GL: uma linguagem específica de domínio para especificação
declarativa e verificável de políticas de degradação gráfica em estações de
controle de VANT}

%----- Bloco de autoria em tres colunas fixas (autores lado a lado).
\author{%
\begin{tabular}{@{}
  >{\centering\arraybackslash}m{0.30\textwidth}@{\hspace{2pt}}
  >{\centering\arraybackslash}m{0.36\textwidth}@{\hspace{2pt}}
  >{\centering\arraybackslash}m{0.30\textwidth}@{}}
Caroline Silva Grizante &
Emilly Emanuely Ribeiro dos Santos &
Marcela Lovatto \\[2pt]
\textit{Engenharia de Computação} &
\textit{Engenharia de Computação} &
\textit{Engenharia de Computação} \\
\textit{Centro Universitário da} &
\textit{Centro Universitário da} &
\textit{Centro Universitário da} \\
\textit{Fundação Hermínio Ometto} &
\textit{Fundação Hermínio Ometto} &
\textit{Fundação Hermínio Ometto} \\
Araras, SP, Brasil &
Araras, SP, Brasil &
Araras, SP, Brasil \\[1pt]
{\footnotesize carolinegrizante105@alunos.fho.edu.br} &
{\footnotesize emillyribeiro@alunos.fho.edu.br} &
{\footnotesize marcela.lovatto@alunos.fho.edu.br}
\end{tabular}
}

\maketitle

%=======================================================================
\begin{abstract}
Estações de controle de solo de veículos aéreos não tripulados exibem
superfícies gráficas de criticidade desigual. Quando o orçamento por quadro é
excedido, decidir qual superfície degradar, e em que grau, é uma decisão de
segurança, em geral codificada em linguagens de propósito geral, sem
verificação automática. Este trabalho apresenta a DRESS-GL, linguagem
específica de domínio externa e não Turing-completa, que expressa essas
políticas no gabarito Dado-Quando-Então, com pontos de junção sobre os passos
de renderização; sem laços nem recursão, cada programa termina por construção.
O processador, escrito em Python sem geradores de analisador, combina análise
léxica por autômato finito determinístico, análise descendente recursiva sobre
gramática cuja propriedade LL(1) é verificada automaticamente e oito
verificações semânticas derivadas do modelo do hospedeiro, e compila as regras
em uma representação intermediária executada por quadro. As políticas
executam, pelo mesmo contrato, em um hospedeiro de custo modelado e no
AgroView-SC, que renderiza em CPU, alimentados por telemetria gravada no
CoppeliaSim. A avaliação examina um cenário excepcional, bateria em reserva
sobre superfície imprópria para pouso, em sete condições. No hospedeiro
simulado, a política manteve todos os quadros no orçamento de 16,67\,ms,
enquanto uma variante sem compensações o excedeu em 13,33\,\% dos quadros. No
AgroView-SC, com relógio por carga calibrado, nenhum programa o excedeu (pior
quadro de 14,33\,ms): a vantagem das compensações depende do modelo de custo.
Com erro de posição de até 0,1\,m, o período mínimo entre reconfigurações
reduziu as transições espúrias em 11\,\% a 22\,\%, sem eliminá-las.
\end{abstract}

\begin{IEEEkeywords}
linguagens específicas de domínio, autômatos finitos, gramáticas livres de
contexto, análise sintática LL(1), análise semântica, sistemas
auto-adaptativos, veículos aéreos não tripulados
\end{IEEEkeywords}

%=======================================================================
\section{Introdução}
\label{sec:introducao}

\subsection{Contexto e domínio}
\label{sec:dominio}

O domínio deste trabalho é a \textit{especificação de políticas de degradação
gráfica adaptativa em estações de controle de solo} (\textit{Ground Control
Stations}, GCS) \textit{de veículos aéreos não tripulados} (VANT)
\textit{de uso agrícola}.

Uma GCS compõe, em um único dispositivo de exibição, superfícies gráficas
heterogêneas: a visão em primeira pessoa transmitida pela aeronave, o mapa de
missão com a rota de pulverização, camadas agronômicas derivadas de índices de
vegetação e um painel de telemetria. Essas superfícies possuem
\textbf{criticidade operacional desigual}. A perda momentânea de atualização da
camada agronômica é irrelevante para a segurança do voo; a mesma perda na visão
em primeira pessoa, durante uma manobra manual, pode ser determinante. A
criticidade tampouco é constante: um mapa de aptidão de pouso é acessório em
voo nominal e torna-se a informação mais importante da tela quando a bateria
entra em reserva.

A composição opera sob \textbf{orçamento temporal rígido}: 16,67\,ms por quadro
para 60\,Hz. Quando a carga excede o orçamento, algo precisa ser sacrificado:
resolução de renderização, frequência de atualização, densidade de geometria ou
a própria presença de uma superfície. Decidir \textit{o que} sacrificar,
\textit{em que ordem}, \textit{sob quais condições de telemetria} e \textit{com
quais garantias mínimas preservadas} é o objeto deste domínio.

Trata-se de uma decisão simultaneamente de engenharia gráfica, de agronomia de
precisão e de segurança operacional. Os especialistas que detêm o conhecimento
necessário para formulá-la, como o agrônomo que estabelece a relevância das
camadas e o piloto que conhece o procedimento de emergência, não são, em geral,
programadores de sistemas gráficos. Essa separação entre quem detém o
conhecimento do domínio e quem detém o conhecimento da implementação é a
condição que van Deursen \textit{et al.}~\cite{vandeursen2000} identificam como
favorável à adoção de uma linguagem dedicada.

O sistema hospedeiro-alvo é o \textbf{AgroView-SC}, compositor gráfico de GCS
desenvolvido pelos autores. O processador não depende do código do compositor:
conhece-o apenas por um modelo declarativo (\texttt{modelo/agroview.toml}), que
descreve superfícies, passos de renderização, instâncias e telemetria, e por
um contrato de hospedeiro, que fixa a troca de planos e leituras a cada quadro.
Neste trabalho, as políticas são executadas em dois hospedeiros que implementam
esse contrato: um hospedeiro simulado, de custo modelado
(Seção~\ref{sec:hospedeiro}), e uma versão reduzida do AgroView-SC que
renderiza em CPU e mede o custo de cada quadro (Seção~\ref{sec:agroview}),
alimentada por telemetria gravada no simulador físico CoppeliaSim
(Seção~\ref{sec:ponte}).

\subsection{Por que uma DSL e não uma linguagem de propósito geral}
\label{sec:relevancia}

Na versão anterior do AgroView-SC, descartada na reimplementação, as políticas
de degradação eram funções em Python registradas em um motor de adaptação. A
Listagem~\ref{lst:python} reconstrói, nesse estilo, a política acionada por
bateria em reserva; a Listagem~\ref{lst:exemplo}, na
Seção~\ref{sec:gramatica}, expressa a mesma intenção em DRESS-GL.

\begin{blocolst}
\begin{lstlisting}[language=Python]
def r_setdown(v, w):
    w["show_land"]   = 1
    w["land_target"] = best_landing_site(w)
    v["MAP"].mode       = 2
    v["AGRO"].target_hz = 2.0
    v["FPV"].crop_ratio = 0.25

E.add(Rule("SafeSetdownAssist", 90,
           lambda t, w: t.battery < 18.0, r_setdown))
\end{lstlisting}
\legendalst{lst:python}{Política de degradação no estilo imperativo da versão
anterior do AgroView-SC (reconstrução ilustrativa).}
\end{blocolst}

Essa formulação apresenta seis deficiências que \textbf{não são defeitos de
implementação, mas consequências de se empregar uma linguagem de propósito
geral} em um domínio com requisitos de verificação.

\textbf{(i) Ausência de verificação de invariantes de segurança.} Nada impede
que uma política escreva \texttt{v["FPV"].enabled = False}, desativando a
superfície de criticidade mais alta. A construção é válida em Python e o erro
só se manifesta em operação. Uma linguagem que conhece a criticidade das
superfícies pode \textbf{rejeitar essa construção em tempo de compilação}; é o
padrão AVOPT de Mernik \textit{et al.}~\cite{mernik2005}.

\textbf{(ii) Ausência de verificação dimensional.} As atribuições
\texttt{target\_hz = 2.0} e \texttt{scale = 2.0} são indistinguíveis para o
interpretador: ambas são \texttt{float}. O domínio, entretanto, distingue
hertz, milissegundos, metros, metros por segundo e fração adimensional, e a
segunda atribuição está fora do intervalo válido de escala. Com unidades
incorporadas à sintaxe, a incompatibilidade é detectável estaticamente.

\textbf{(iii) Ordem de avaliação relevante e não declarada.} Duas políticas
simultaneamente ativas podem escrever o mesmo atributo com valores distintos; o
resultado depende da ordem de registro no código-fonte, que não faz parte da
especificação. Uma linguagem dedicada pode impor \textbf{composição
monotônica} por atributo, com mínimo para atributos numéricos e conjunção para
booleanos, operações comutativas e associativas e, portanto, independentes da
ordem. Para atributos enumerados, sem ordem natural, a escrita é resolvida pela
prioridade declarada e o conflito é registrado.

\textbf{(iv) Não terminação e custo ilimitado.} A política é código
Turing-completo executado dentro do laço de renderização. Um laço mal
formulado, uma chamada custosa ou uma recursão inadvertida violam o orçamento.
Na versão anterior, a chamada correspondente a \texttt{best\_landing\_site()}
percorria uma matriz de $512 \times 512$ elementos a cada quadro em que a
política estava ativa, sem limite declarado; esse custo não foi medido, pois o
código foi descartado. Para um programa arbitrário em linguagem
Turing-completa, decidir se a execução termina é indecidível e, \textit{a
fortiori}, não há limite superior geral para seu custo~\cite{hopcroft2007}. Uma
linguagem sem laços, sem recursão e com conjunto fechado de ações termina
\textit{por construção}: a cada quadro, cada condição é avaliada uma vez,
visitando cada nó de sua árvore, e cada efeito é composto uma vez, de modo que
o trabalho do motor é limitado por um polinômio no tamanho do programa,
conhecido após a compilação. O custo de executar cada ação fica com o
hospedeiro e é declarado no seu modelo.

\textbf{(v) Inauditabilidade.} A Listagem~\ref{lst:python} é de difícil leitura
para o agrônomo e para o piloto. O gabarito Dado-Quando-Então aproxima a
especificação do vocabulário do domínio; a revisão de~\cite{rodrigues2017}
mostra, contudo, que esse ganho precisa ser medido com participantes, e não
presumido pelos projetistas.

\textbf{(vi) Espaço de configuração ilimitado.} Como uma função arbitrária pode
produzir qualquer valor, o conjunto de configurações alcançáveis não é
enumerável. Na DRESS-GL, o plano de cada quadro é função apenas do conjunto de
regras ativas; um programa com $n$ regras alcança, portanto, no máximo $2^n$
planos, cada um verificável por inspeção.

Bibliotecas e padrões de projeto sobre a linguagem hospedeira atenuam (v), mas
não afetam (i)--(iv), pois não é possível restringir o poder expressivo de uma
linguagem de propósito geral sem deixar de usá-la. Arquivos JSON ou YAML validam
estrutura, não significado, e não expressam condição, ponto de junção nem
contrato. É a situação em que~\cite{vandeursen2000} e~\cite{mernik2005}
recomendam uma linguagem dedicada. O principal risco apontado
por~\cite{vandeursen2000}, o custo de projeto e manutenção da linguagem,
justifica-se pela recorrência do artefato: as políticas são reescritas a cada
nova aeronave, cultura, regulamentação e perfil de missão, por pessoas
distintas das que implementaram o compositor.

\subsection{Contribuições e organização}
\label{sec:contrib}

As contribuições deste trabalho são: (i) o projeto da DRESS-GL, cuja restrição
deliberada de poder expressivo garante terminação e trabalho limitado por
quadro; (ii) um processador completo, com analisadores escritos manualmente,
verificação automática da propriedade LL(1), tabela de símbolos derivada do
modelo do hospedeiro e oito verificações semânticas; (iii) um contrato de
hospedeiro com duas implementações, uma de custo modelado e outra de custo
medido, e uma ponte de telemetria com o CoppeliaSim baseada em gravações
reproduzíveis; e (iv) uma avaliação experimental reprodutível de um cenário
excepcional em sete condições. O código-fonte, os testes, o registro numerado
de decisões de projeto (\texttt{docs/decisions.md}; ao longo do texto, D$n$
indica a decisão $n$), as instruções de instalação (\texttt{README.md}) e o
manual da linguagem (\texttt{docs/manual.md}) estão disponíveis no repositório
\url{https://github.com/EmillyeRibeiro/dressgl}, na \textit{tag}
\texttt{v3.0}.

A Seção~\ref{sec:fundamentacao} apresenta a fundamentação teórica e os
trabalhos relacionados; a Seção~\ref{sec:desenvolvimento}, os requisitos, a
gramática, as ferramentas e a arquitetura do processador e dos hospedeiros; a
Seção~\ref{sec:resultados-discussao}, os resultados de compilação, de teste e
de execução; e a Seção~\ref{sec:conclusao}, as conclusões.

%=======================================================================
\section{Fundamentação Teórica}
\label{sec:fundamentacao}

\subsection{Linguagens específicas de domínio}
\label{sec:dsl}

Van Deursen, Klint e Visser~\cite{vandeursen2000} definem uma DSL como uma
linguagem de programação ou de especificação executável que oferece, por meio
de notações e abstrações apropriadas, poder expressivo focado em, e usualmente
restrito a, um domínio de problema particular. DSLs são tipicamente
declarativas e podem ser vistas simultaneamente como linguagens de
especificação e de programação. Entre os benefícios estão expressividade,
autodocumentação, validação e otimização no nível do domínio e envolvimento do
especialista; entre os riscos, custo de projeto, implementação e manutenção,
custo de treinamento, dificuldade de delimitar escopo e eventual perda de
eficiência.

Mernik, Heering e Sloane~\cite{mernik2005} organizam o desenvolvimento de DSLs
em padrões de decisão, análise, projeto e implementação. Dois padrões de
decisão aplicam-se aqui: \textbf{notação}, quando a notação do domínio é mal
suportada por linguagens de propósito geral, e \textbf{AVOPT}
(\textit{analysis, verification, optimization, parallelization,
transformation}), quando análises específicas do domínio são desejáveis, mas
difíceis ou impossíveis sobre uma linguagem Turing-completa. Quanto ao projeto,
este trabalho segue o padrão de \textit{invenção de linguagem}: uma linguagem
externa sem relação com uma linguagem hospedeira; quanto à implementação, o
padrão \textit{interpretador}.

Nascimento \textit{et al.}~\cite{nascimento2012} mapeiam domínios, ferramentas e
processos de desenvolvimento de DSLs e discutem as \textbf{linguagens
específicas de domínio orientadas a aspectos} (DSAL), que combinam DSLs com
programação orientada a aspectos para modularizar interesses transversais. A
cláusula \texttt{WHEN} da DRESS-GL, que declara pontos de junção
(\texttt{BEFORE}, \texttt{AFTER}, \texttt{AROUND}) sobre passos de renderização,
insere a linguagem nessa categoria: a política de degradação atravessa todos os
passos e não pode ser modularizada no interior de nenhum deles.

Rodrigues, Campos e Zorzo~\cite{rodrigues2017} constatam que a avaliação de
usabilidade de DSLs é frequentemente omitida ou subjetiva. Como a legibilidade
por especialistas é um objetivo da DRESS-GL, essa constatação é tratada como
requisito: sustentar o objetivo exige um estudo com participantes, não
realizado neste trabalho (Seção~\ref{sec:ameacas}).

Boukham, Younsi Dahbi e Chiadmi~\cite{boukham2025} revisam DSLs para
processamento de grafos e observam que a separação entre a abstração do domínio
e as plataformas-alvo tende a ser resolvida por representação intermediária. Na
DRESS-GL, o alvo é também um grafo, o de passos de renderização, e a separação
é feita pela representação intermediária de condições e efeitos
(Seção~\ref{sec:backend}); é essa separação que permite executar o mesmo
programa compilado em dois hospedeiros distintos. Breuker~\cite{breuker2014}
analisa DSLs para aprendizado de máquina sob a perspectiva da engenharia
dirigida a modelos, evidenciando o valor de transformar especificações
declarativas em artefatos executáveis, como faz o processador entre o arquivo
\texttt{.dgl} e o motor.

\subsection{Linguagens regulares e análise léxica}
\label{sec:lexica}

A estrutura lexical de uma linguagem de programação é tipicamente regular:
cada classe de \textit{token} é descrita por uma expressão regular, e a união
dessas expressões é reconhecida por um autômato finito determinístico
(AFD)~\cite{hopcroft2007}. O analisador léxico simula o AFD sobre a entrada e,
quando há mais de um prefixo aceito, aplica a regra do lexema mais longo
(\textit{maximal munch}); palavras reservadas são reconhecidas como
identificadores e reclassificadas por consulta a uma tabela, o que mantém o
autômato pequeno~\cite{aho2007}.

\subsection{Gramáticas livres de contexto e análise LL(1)}
\label{sec:ll1}

A estrutura sintática é descrita por uma gramática livre de contexto. Um
analisador descendente recursivo associa uma função a cada não terminal e
escolhe a produção a expandir pelo \textit{token} corrente. Isso dispensa
retrocesso quando a gramática é LL(1): para cada não terminal $A$, os conjuntos
das suas produções $A \to \alpha$, com $F = \mathrm{FIRST}(\alpha)$,
\begin{equation*}
\mathrm{SELECT}(A \to \alpha) =
\begin{cases}
F, & \varepsilon \notin F\\
(F \setminus \{\varepsilon\}) \cup \mathrm{FOLLOW}(A), & \varepsilon \in F
\end{cases}
\end{equation*}
devem ser disjuntos dois a dois~\cite{aho2007}. Recursão à esquerda e prefixos
comuns violam essa condição e são removidos por eliminação de recursão e
fatoração à esquerda. Para recuperação de erros, o modo pânico descarta
\textit{tokens} até encontrar um elemento de um conjunto de sincronização,
permitindo relatar vários erros em uma única execução~\cite{aho2007}.

\subsection{Análise semântica e tabela de símbolos}
\label{sec:semantica-teoria}

Restrições que não são expressáveis por gramática livre de contexto, como a
concordância entre uso e declaração de um nome, são deslocadas para a análise
semântica, que percorre a árvore sintática abstrata (AST) consultando uma
tabela de símbolos~\cite{aho2007}. Na DRESS-GL, essa fase concentra também as
restrições do domínio: compatibilidade de unidades, pertinência a domínios de
valores e invariantes de criticidade.

\subsection{Trabalhos relacionados}
\label{sec:relacionados}

A DRESS-ML, de Alves \textit{et al.}~\cite{alves2022}, é uma DSL para modelagem
de cenários excepcionais e comportamentos auto-adaptativos em aplicações
baseadas em drones. Combina o gabarito Dado-Quando-Então, oriundo do
desenvolvimento guiado por comportamento, com conceitos de programação
orientada a aspectos; é implementada sobre o JetBrains MPS e gera, por
transformação de modelo para texto, envoltórios executáveis para o simulador
Dragonfly. A Stitch, de Cheng, Garlan e Schmerl~\cite{cheng2012}, representa
estratégias de adaptação baseada em arquitetura no arcabouço Rainbow, com
operadores, táticas e estratégias selecionadas por teoria da utilidade. A
Tabela~\ref{tab:critica} relaciona essas e outras abordagens aplicáveis ao
domínio e a limitação central de cada uma.

\begin{table*}[t]
\caption{Abordagens existentes e suas limitações para o domínio}
\label{tab:critica}
\centering
\scriptsize
\setlength{\tabcolsep}{3pt}
\renewcommand{\arraystretch}{0.95}
\begin{tabularx}{\textwidth}{@{}>{\raggedright\arraybackslash}p{0.17\textwidth}
                               >{\raggedright\arraybackslash}p{0.19\textwidth}L@{}}
\toprule
\textbf{Abordagem} & \textbf{Natureza} & \textbf{Limitação central para o domínio} \\
\midrule
Funções e expressões lambda em linguagem de propósito geral (versão anterior do
AgroView-SC) &
Código Turing-completo embutido no laço de renderização &
Sem verificação de invariantes de criticidade nem checagem dimensional;
terminação não garantida e custo ilimitado; espaço de configuração não
enumerável; de difícil auditoria pelo especialista \\
\addlinespace
Arquivos de configuração em JSON, YAML ou INI &
Formato de dados sem semântica de domínio &
Não expressam condição, ponto de junção nem contrato; a validação disponível é
de conformidade a esquema estrutural, não de significado \\
\addlinespace
DRESS-ML~\cite{alves2022} &
DSL de adaptação para drones; Dado-Quando-Então e aspectos, sobre JetBrains MPS &
\textbf{Referência mais próxima.} Modela aeronave, missão, sensores e atuadores,
mas não o subsistema de visualização; não possui atuadores gráficos nem noção
de orçamento por quadro; validação relatada como prova de conceito em simulador
único, e os autores registram a ausência de estudos empíricos de usabilidade \\
\addlinespace
Stitch~\cite{cheng2012} &
DSL de táticas e estratégias de adaptação arquitetural, no arcabouço Rainbow &
Granularidade de componente arquitetural e escala de tempo de segundos a
minutos; inadequada a decisões por quadro em janela de 16,67\,ms; não modela
recursos gráficos \\
\addlinespace
Linguagens de sombreamento (GLSL, HLSL) e grafos de material &
Especificação de aparência visual &
Especificam a aparência de cada fragmento, não a \textit{política} de alocação
do orçamento entre superfícies de criticidade distinta \\
\addlinespace
Escalabilidade de motores gráficos: perfis de qualidade e escala dinâmica de
resolução &
Configuração e heurística embutidas no motor &
Operam sobre um eixo único de fidelidade global e supõem critério perceptual
homogêneo; desconhecem criticidade desigual e variável no tempo \\
\addlinespace
Linguagens de plano de missão e árvores de comportamento &
Descrição de missão e de controle &
Adaptam o comportamento da aeronave, não os recursos de visualização da
estação de controle \\
\addlinespace
Linguagens de descrição de cenário para simulação &
Definição declarativa de cenário &
Descrevem os estímulos do ambiente, não a resposta adaptativa do sistema de
exibição \\
\addlinespace
Definição declarativa de camadas em aviônica embarcada &
Superfícies e prioridades definidas em tempo de projeto &
Prioridades fixadas antes da execução; sem reconfiguração dependente do estado
corrente da missão \\
\bottomrule
\end{tabularx}
\end{table*}

Nenhuma das abordagens reúne, simultaneamente: (a) notação declarativa voltada
ao especialista; (b) pontos de junção sobre o grafo de passos de renderização;
(c) verificação estática de invariantes de criticidade; (d) restrição sintática
que garanta terminação; e (e) contrato temporal explícito. A DRESS-GL herda da
DRESS-ML o gabarito Dado-Quando-Então e a organização orientada a aspectos, e
dela se distingue em três pontos: adapta o subsistema de visualização da
estação, e não o comportamento da aeronave; submete as políticas a um
orçamento temporal com contratos declarados em \texttt{ASSERT}; e é
implementada como linguagem externa com analisadores próprios, e não sobre um
ambiente de modelagem, o que expõe as fases de tradução e permite verificações
específicas do domínio.

%=======================================================================
\section{Desenvolvimento}
\label{sec:desenvolvimento}

\subsection{Requisitos e escopo}
\label{sec:requisitos}

\textbf{Objetivo geral.} Projetar e implementar a DRESS-GL e um processador que
realize as análises léxica, sintática e semântica, compile as políticas
validadas e as execute sobre hospedeiros que implementam o modelo e o contrato
do compositor AgroView-SC, um deles com custo de quadro medido.

\textbf{Funcionalidades da linguagem.}
\begin{itemize}
\item[\textbf{F1}] Declarar a política global (\texttt{POLICY}), com orçamento
de quadro, reserva e período mínimo entre reconfigurações
(\texttt{cooldown}).
\item[\textbf{F2}] Declarar regras nomeadas e priorizadas segundo o gabarito
\texttt{RULE $\dots$ GIVEN $\dots$ WHEN $\dots$ THEN $\dots$ [ASSERT] $\dots$
END}.
\item[\textbf{F3}] Expressar, em \texttt{GIVEN}, predicados sobre telemetria
com operadores relacionais, conectivos \texttt{AND}, \texttt{OR} e
\texttt{NOT}, pertinência a conjuntos enumerados e unidades de medida.
\item[\textbf{F4}] Declarar, em \texttt{WHEN}, o ponto de junção no grafo de
renderização, nas fases \texttt{BEFORE}, \texttt{AFTER} e \texttt{AROUND}.
\item[\textbf{F5}] Expressar, em \texttt{THEN}, um conjunto \textbf{fechado} de
ações: atribuição de atributos de superfície (\texttt{SET}), escalonamento de
instâncias e de resolução (\texttt{SCALE}), injeção e supressão de passos
(\texttt{INJECT}, \texttt{SUPPRESS}), destaque de extremos de campos escalares
(\texttt{HIGHLIGHT}) e marcação de obsolescência (\texttt{MARK}).
\item[\textbf{F6}] Declarar, em \texttt{ASSERT}, contratos sobre o efeito das
ações, verificados estaticamente e em execução.
\item[\textbf{F7}] \textbf{Não} oferecer laços, recursão, definição de funções
nem aritmética: restrição deliberada que garante terminação
(Seção~\ref{sec:relevancia}).
\end{itemize}

\textbf{Requisitos do processador.}
\begin{itemize}
\item[\textbf{P1}] Análise léxica por AFD, com literais numéricos anotados por
unidade e registro de linha e coluna de cada \textit{token}.
\item[\textbf{P2}] Análise sintática descendente recursiva sobre gramática
LL(1), produzindo AST tipada, com recuperação de erros.
\item[\textbf{P3}] Tabela de símbolos construída a partir do modelo do
hospedeiro, com resolução de nomes por espaço de nomes.
\item[\textbf{P4}] Oito verificações semânticas com diagnóstico posicionado:
\textbf{V1} referência indefinida; \textbf{V2} incompatibilidade de unidade;
\textbf{V3} valor fora do domínio admissível; \textbf{V4} violação de
invariante de criticidade; \textbf{V5} ponto de junção inválido ou
ambíguo; \textbf{V6} conflito de escrita entre regras de prioridade igual;
\textbf{V7} regra logicamente inalcançável; \textbf{V8} asserção sem
instrumento de medição.
\item[\textbf{P5}] Compilação de cada regra validada em representação
intermediária (RI) de condição e efeitos, de modo que o hospedeiro aplique um
plano de ações sem executar código arbitrário.
\item[\textbf{P6}] Execução por quadro, com composição monotônica entre regras
ativas e período mínimo entre reconfigurações.
\item[\textbf{P7}] Rastreabilidade entre regra, ações, alvos e asserções.
\item[\textbf{P8}] Contrato de hospedeiro independente da implementação, com
um hospedeiro de custo medido e entrada de telemetria por simulador físico,
sem que o núcleo do processador dependa de bibliotecas externas.
\end{itemize}

A Tabela~\ref{tab:escopo} delimita o projeto, conforme a recomendação
de~\cite{mernik2005} de fechar o escopo antes do projeto da linguagem. O
domínio que motiva a linguagem compreende seis cenários excepcionais: bateria
em reserva sobre superfície imprópria para pouso; perda de sinal de
posicionamento com transição para guiagem manual; vento acima do limite de
deriva de pulverização; degradação do sensor de obstáculos; tanque de calda
esgotado; e a ocorrência simultânea desses eventos. O \textbf{escopo avaliado}
é deliberadamente mais estreito (D57): \textbf{um único cenário}, o primeiro,
examinado em \textbf{sete condições} de execução (Tabela~\ref{tab:condicoes},
Seção~\ref{sec:avaliacao}). Os outros cinco exigem campos de telemetria que o
modelo do hospedeiro ainda não declara e que a cena do simulador não fornece;
ficam como trabalho futuro (Seção~\ref{sec:conclusao}).

\begin{table}[htbp]
\caption{Delimitação do escopo do projeto}
\label{tab:escopo}
\centering
\footnotesize
\begin{tabularx}{\columnwidth}{@{}LL@{}}
\toprule
\textbf{Dentro do escopo} & \textbf{Fora do escopo} \\
\midrule
Linguagem externa e textual, arquivos \texttt{.dgl} &
Editor gráfico, ambiente projecional, sintaxe visual \\
\addlinespace
Analisadores léxico e sintático escritos manualmente &
Geradores de analisador; inferência de tipos; polimorfismo \\
\addlinespace
Oito verificações semânticas (V1--V8) &
Verificação formal por provador; \textit{model checking} \\
\addlinespace
\textit{Back end} interpretador sobre dois hospedeiros: simulado e AgroView-SC
renderizado em CPU &
Geração de código C/C++ certificável; compositor acelerado por GPU \\
\addlinespace
Atuadores gráficos: resolução, frequência, densidade de instâncias, modo de
visualização, injeção e supressão de passo &
Atuadores de rede, de armazenamento ou de controle de voo \\
\addlinespace
Telemetria de um único VANT, de cenário sintético ou gravada no CoppeliaSim &
Enxames; telemetria de aeronave real; laço fechado com o simulador \\
\addlinespace
Um cenário excepcional em sete condições; quatro programas, com três regras
distintas &
Os outros cinco cenários; biblioteca completa de procedimentos de emergência \\
\addlinespace
Custo modelado e custo medido em CPU, com relógio por carga calibrado &
Certificação aeronáutica; ensaio em voo; medição sobre GPU \\
\addlinespace
Reprodução de gravações sem interface gráfica &
Janela de visualização (reservada ao extra \texttt{gui}) \\
\bottomrule
\end{tabularx}
\end{table}

\subsection{Gramática formal}
\label{sec:gramatica}

A Listagem~\ref{lst:ebnf} apresenta a gramática em EBNF. As definições
regulares do topo descrevem as classes de \textit{tokens}; as demais
produções, a sintaxe. A unidade de medida é parte do \textit{token}
\texttt{NUM}, e não da gramática sintática: a linguagem léxica permanece
regular, e cada literal numérico chega à análise semântica já anotado com sua
unidade, o que viabiliza a verificação V2.

\begin{blocolst}
\begin{lstlisting}[basicstyle=\ttfamily\fontsize{6.5}{7.8}\selectfont]
(* ---- lexico ---- *)
letra      = "A".."Z" | "a".."z" | "_" ;
digito     = "0".."9" ;
(* palavra: reservada, BOOL, ENUM ou IDENT *)
palavra    = letra { letra | digito } ;
NUM        = digito { digito } [ "." digito { digito } ]
             [ { " " | "\t" } UNIDADE ] ;
UNIDADE    = "frames" | "m/s" | "ms" | "Hz" | "m" | "%" ;
STRING     = '"' { caractere exceto '"' e quebra
                   de linha } '"' ;
comentario = "#" { caractere exceto quebra de linha } ;

(* ---- sintaxe ---- *)
programa   = { declaracao } ;
declaracao = politica | regra ;
politica   = "POLICY" IDENT "{" { atrib } "}" ;
atrib      = param "=" literal ";" ;
param      = "budget" | "reserve" | "cooldown" ;
regra      = "RULE" IDENT [ "PRIORITY" NUM ]
             "GIVEN" expr "WHEN" juncao "THEN" acoes
             [ "ASSERT" expr ] "END" ;
expr       = conj { "OR" conj } ;
conj       = neg { "AND" neg } ;
neg        = [ "NOT" ] atomo ;
atomo      = "(" expr ")" | referencia resto ;
resto      = op_rel literal | "IN" conjunto ;
op_rel     = "<" | "<=" | ">" | ">=" | "==" | "!=" ;
conjunto   = "{" ENUM { "," ENUM } "}" ;
referencia = IDENT [ "." IDENT ] | alvo "." IDENT ;
alvo       = "view" STRING | "pass" STRING ;
juncao     = fase "pass" STRING ;
fase       = "BEFORE" | "AFTER" | "AROUND" ;
acoes      = acao ";" { acao ";" } ;
acao       = "SET"       alvo "." IDENT "=" literal
           | "SCALE"     recurso STRING "BY" NUM
           | "INJECT"    "pass" STRING fase_rel STRING
           | "SUPPRESS"  "pass" STRING
           | "HIGHLIGHT" extremo referencia
                         [ "WITHIN" literal ]
           | "MARK"      alvo "STALE" ;
recurso    = "instances" | "resolution" ;
fase_rel   = "BEFORE" | "AFTER" ;
extremo    = "ARGMAX" | "ARGMIN" ;
literal    = NUM | ENUM | BOOL | STRING ;
\end{lstlisting}
\legendalst{lst:ebnf}{Gramática da DRESS-GL em EBNF: definições regulares dos
\textit{tokens} (topo) e produções sintáticas.}
\end{blocolst}

Uma palavra é classificada, nesta ordem, como palavra reservada (33 no total),
booleano (\texttt{true}, \texttt{false}), enumerado, se começa por maiúscula e
contém apenas maiúsculas, dígitos e sublinhado (\texttt{WATER},
\texttt{STANDING\_CROP}), ou identificador. Por consequência, nomes de regra e
de política são escritos em PascalCase (\texttt{SafeSetdownAssist}); um nome
como \texttt{GPS\_LOSS} seria lido como enumerado. A exigência de que
\texttt{PRIORITY} receba um inteiro sem unidade é verificada pelo analisador
sintático, que registra o erro sem abandonar a regra.

Para a verificação, as construções EBNF foram reescritas em BNF (repetições e
opcionais viram não terminais com alternativa vazia), resultando em 33 não
terminais e 52 terminais (palavras reservadas, sete operadores, sete sinais de
pontuação e as classes \texttt{IDENT}, \texttt{NUM}, \texttt{STRING},
\texttt{ENUM} e \texttt{BOOL}). O módulo \texttt{gramatica\_ll1} calcula FIRST e
FOLLOW por ponto fixo, deriva os conjuntos SELECT e verifica que eles são
disjuntos para todo não terminal; um teste automatizado falha se surgir
conflito LL(1) ou não terminal inalcançável, e outro confirma que os terminais
da gramática coincidem com os \textit{tokens} do analisador léxico. A
Tabela~\ref{tab:select} mostra os casos decisivos: os não terminais anuláveis,
em que a escolha depende de FOLLOW, e os que têm alternativas não vazias. Em
\texttt{atomo}, a fatoração à esquerda em \texttt{resto} elimina o prefixo
comum \texttt{referencia}; em \texttt{campo}, o ponto não pertence a
FOLLOW(\texttt{referencia}), o que permite referências simples, como
\texttt{frame\_budget}, sem perder a propriedade LL(1). Os demais não terminais
com alternativas (\texttt{declaracao}, \texttt{param}, \texttt{op\_rel},
\texttt{alvo}, \texttt{fase}, \texttt{acao}, \texttt{recurso},
\texttt{fase\_rel}, \texttt{extremo}, \texttt{literal}) começam cada
alternativa por um terminal distinto.

\begin{table}[htbp]
\caption{Conjuntos SELECT dos não terminais anuláveis e dos que exigem
fatoração}
\label{tab:select}
\centering
\scriptsize
\setlength{\tabcolsep}{3pt}
\begin{tabularx}{\columnwidth}{@{}l >{\raggedright\arraybackslash}p{0.30\columnwidth} L@{}}
\toprule
\textbf{Não terminal} & \textbf{Alternativa} & \textbf{SELECT} \\
\midrule
\texttt{decls}      & \texttt{declaracao decls} & \texttt{POLICY}, \texttt{RULE} \\
                    & $\varepsilon$ & \texttt{\$} \\
\texttt{atribs}     & \texttt{atrib atribs} & \texttt{budget}, \texttt{reserve}, \texttt{cooldown} \\
                    & $\varepsilon$ & \texttt{\}} \\
\texttt{prio}       & \texttt{PRIORITY NUM} & \texttt{PRIORITY} \\
                    & $\varepsilon$ & \texttt{GIVEN} \\
\texttt{assert\_opt}& \texttt{ASSERT expr} & \texttt{ASSERT} \\
                    & $\varepsilon$ & \texttt{END} \\
\texttt{expr\_r}    & \texttt{OR conj expr\_r} & \texttt{OR} \\
                    & $\varepsilon$ & \texttt{WHEN}, \texttt{END}, \texttt{)} \\
\texttt{conj\_r}    & \texttt{AND neg conj\_r} & \texttt{AND} \\
                    & $\varepsilon$ & \texttt{OR}, \texttt{WHEN}, \texttt{END}, \texttt{)} \\
\texttt{campo}      & \texttt{. IDENT} & \texttt{.} \\
                    & $\varepsilon$ & \texttt{< <= > >= == !=}, \texttt{IN}, \texttt{WITHIN}, \texttt{;} \\
\texttt{enums}      & \texttt{, ENUM enums} & \texttt{,} \\
                    & $\varepsilon$ & \texttt{\}} \\
\texttt{acoes\_r}   & \texttt{acao ; acoes\_r} & \texttt{SET}, \texttt{SCALE}, \texttt{INJECT}, \texttt{SUPPRESS}, \texttt{HIGHLIGHT}, \texttt{MARK} \\
                    & $\varepsilon$ & \texttt{ASSERT}, \texttt{END} \\
\texttt{within}     & \texttt{WITHIN literal} & \texttt{WITHIN} \\
                    & $\varepsilon$ & \texttt{;} \\
\addlinespace
\texttt{neg}        & \texttt{NOT atomo} & \texttt{NOT} \\
                    & \texttt{atomo} & \texttt{(}, \texttt{IDENT}, \texttt{view}, \texttt{pass} \\
\texttt{atomo}      & \texttt{( expr )} & \texttt{(} \\
                    & \texttt{referencia resto} & \texttt{IDENT}, \texttt{view}, \texttt{pass} \\
\texttt{referencia} & \texttt{IDENT campo} & \texttt{IDENT} \\
                    & \texttt{alvo . IDENT} & \texttt{view}, \texttt{pass} \\
\texttt{resto}      & \texttt{op\_rel literal} & \texttt{< <= > >= == !=} \\
                    & \texttt{IN conjunto} & \texttt{IN} \\
\bottomrule
\end{tabularx}
\end{table}

A Listagem~\ref{lst:exemplo} exibe um programa completo, correspondente ao
cenário de bateria em reserva sobre superfície imprópria para pouso. Em relação
à Listagem~\ref{lst:python}, ele declara explicitamente a condição, o ponto de
junção, as ações e o contrato a preservar.

\begin{blocolst}
\begin{lstlisting}[style=dsl]
POLICY AgroViewSC {
  budget   = 16.67 ms;
  reserve  =  2.00 ms;
  cooldown =    10 frames;
}

RULE SafeSetdownAssist PRIORITY 90
  GIVEN drone.battery < 18 %
    AND terrain.below IN { WATER, STANDING_CROP }
  WHEN  AROUND pass "TerrainPass"
  THEN  INJECT pass "SuitabilityPass" AFTER "TerrainPass";
        SET view "MAP".mode = SUITABILITY;
        HIGHLIGHT ARGMAX terrain.suitability WITHIN 400 m;
        SET view "AGRO".hz = 2 Hz;
        SCALE instances "crops" BY 0.25;
  ASSERT frame_budget <= 16.67 ms
     AND view "FPV".scale >= 0.75
END
\end{lstlisting}
\legendalst{lst:exemplo}{Política de assistência a pouso seguro em DRESS-GL
(\texttt{examples/\allowbreak safe\_setdown.dgl}).}
\end{blocolst}

\subsection{Ferramentas e estratégia de implementação}
\label{sec:ferramentas}

O processador foi implementado em Python~3.11. O núcleo (pacote
\texttt{dressgl}: compilador, motor, contrato e hospedeiro simulado) usa apenas
a biblioteca padrão: \texttt{tomllib} para o modelo do hospedeiro e os
cenários, \texttt{dataclasses} imutáveis para \textit{tokens}, AST e RI, e
\texttt{argparse} para as interfaces de linha de comando. Nenhum gerador de
analisador (PLY, ANTLR, Flex/Bison) foi empregado: a implementação manual expõe
a simulação do AFD, a correspondência entre funções e não terminais, a
recuperação de erros e a integração entre a tabela de símbolos e as
verificações do domínio, aspectos que um gerador encapsula. As demais
dependências são opcionais e declaradas como extras do
\texttt{pyproject.toml}: NumPy para a renderização do AgroView-SC
(\texttt{render}), o cliente da API remota ZeroMQ do CoppeliaSim
(\texttt{coppelia}) e \texttt{pytest}~9.1.1 com NumPy para os testes
(\texttt{dev}, D59). Um teste executa o núcleo em subprocesso sem NumPy
disponível e confirma que ele não o importa. As versões exatas estão em
\texttt{results/requirements-lock.txt}, e o ambiente de execução (Windows~11,
compilação 26200, processador Intel Core i7-13620H com 16 \textit{threads}), em
\texttt{results/ambiente.json}.

Quanto à estratégia, a DRESS-GL é \textbf{compilada para uma RI e
interpretada}: cada regra validada é traduzida em uma condição, avaliada por
caminhamento da AST da expressão, e em uma tupla imutável de efeitos, que o
motor compõe a cada quadro em um plano de ações. Não há geração de código-alvo.
As decisões de projeto estão registradas, numeradas, em
\texttt{docs/decisions.md}, e o manual do usuário, em \texttt{docs/manual.md}.

\subsection{Arquitetura do processador e dos hospedeiros}
\label{sec:arquitetura}

A Fig.~\ref{fig:arquitetura} mostra o fluxo de dados. A interrupção entre fases
é deliberada: erros léxicos ou sintáticos impedem a análise semântica, para
evitar diagnósticos em cascata sobre uma AST parcial, e erros semânticos
impedem a geração da RI; avisos não bloqueiam. Dentro de cada fase, porém,
nenhum erro interrompe a análise: todos são relatados de uma vez.

\begin{figure*}[t]
\centering
\begin{adjustbox}{max width=\textwidth}
\begin{tikzpicture}[node distance=5mm, >={Stealth[length=1.6mm]},
  b/.style={draw, rounded corners=2pt, align=center, font=\scriptsize,
            minimum height=9mm, inner sep=3pt},
  l/.style={font=\tiny}]
\node[b] (src) {Fonte\\\texttt{.dgl}};
\node[b, right=of src] (lex) {Léxico\\(AFD)};
\node[b, right=of lex] (sin) {Sintático\\(desc.\ recursivo)};
\node[b, right=of sin] (sem) {Semântico\\(V1--V8, C1, C2)};
\node[b, right=of sem] (ger) {Gerador\\de RI};
\node[b, right=of ger] (mot) {Motor\\(\texttt{cooldown}, composição,\\\texttt{ASSERT})};
\node[b, below=10mm of sem] (mod) {Modelo do hospedeiro\\\texttt{agroview.toml}};
\node[b, below=10mm of mot] (con) {Contrato do hospedeiro\\(\texttt{hospedeiro.py}, D40)};
\node[b, right=6mm of con] (tr) {\textit{Trace} JSONL\\e métricas};
\node[b, below=10mm of con, xshift=-23mm] (sim) {Hospedeiro simulado\\(custo modelado,\\Eq.~(\ref{eq:custo}))};
\node[b, below=10mm of con, xshift=23mm] (agv) {AgroView-SC\\(render em CPU,\\relógio \texttt{tempo}/\texttt{carga})};
\node[b, below=27mm of con] (cop) {CoppeliaSim 4.10\\(ponte ZeroMQ,\\gravações JSONL)};
\draw[->] (src) -- (lex);
\draw[->] (lex) -- node[above, l] {tokens} (sin);
\draw[->] (sin) -- node[above, l] {AST} (sem);
\draw[->] (sem) -- node[above, l] {AST} (ger);
\draw[->] (ger) -- node[above, l] {RI} (mot);
\draw[->] (mod) -- node[left, l] {símbolos} (sem);
\draw[->] (mod) -- node[above, l] {custos, padrões} (con);
\draw[->] ([xshift=-5mm]mot.south) -- node[left, l] {plano $q$} ([xshift=-5mm]con.north);
\draw[->] ([xshift=5mm]con.north) -- node[right, l] {leituras $q{+}1$} ([xshift=5mm]mot.south);
\draw[->] (mot.east) -| (tr.north);
\draw[<->] (con) -- (sim);
\draw[<->] (con) -- (agv);
\draw[->, dashed] (cop) -- node[sloped, below, l] {entradas} (sim);
\draw[->, dashed] (cop) -- node[sloped, below, l] {entradas} (agv);
\end{tikzpicture}
\end{adjustbox}
\caption{Arquitetura do processador DRESS-GL, dos dois hospedeiros e do laço
de execução por quadro. As linhas tracejadas indicam telemetria gravada no
CoppeliaSim e reproduzida pelos hospedeiros.}
\label{fig:arquitetura}
\end{figure*}

\subsubsection{Analisador léxico}
\label{sec:lexer}
O módulo \texttt{lexer} simula diretamente o AFD da Fig.~\ref{fig:afd}: o
primeiro caractere escolhe o ramo, e um laço consome o maior lexema aceito.
Depois de reconhecer um número, o analisador examina, sem consumir, se após
espaços ou tabulações segue uma unidade não continuada por letra ou dígito;
apenas nesse caso a unidade é anexada ao \textit{token}. Assim,
\texttt{10 frames} produz um único \texttt{NUM} com unidade \texttt{frames},
enquanto \texttt{10 mode} produz \texttt{NUM} seguido de \texttt{IDENT}. As
unidades são testadas da mais longa para a mais curta (\texttt{m/s} e
\texttt{ms} antes de \texttt{m}). Um caractere sem transição ou uma cadeia não
terminada geram um \textit{token} \texttt{ERRO} com linha e coluna, e a
varredura continua. Identificadores aceitam apenas letras ASCII, limitação
documentada e coberta por teste.

\begin{figure}[htbp]
\centering
\begin{adjustbox}{max width=\columnwidth}
\begin{tikzpicture}[>={Stealth[length=1.5mm]}, shorten >=1pt,
  every state/.style={draw, circle, minimum size=7mm, inner sep=0pt,
                      font=\scriptsize},
  font=\tiny]
\node[state, initial, initial text={}] (q0) at (0,0)      {$q_0$};
\node[state, accepting] (nid)  at (2.2,3.4)  {ID};
\node[state, accepting] (nint) at (2.2,2.0)  {INT};
\node[state]            (npt)  at (3.6,2.0)  {$q_p$};
\node[state, accepting] (nfr)  at (5.0,2.0)  {FRA};
\node[state]            (ns)   at (2.2,0.4)  {$q_s$};
\node[state, accepting] (nsf)  at (3.6,0.4)  {STR};
\node[state, accepting] (no1)  at (2.2,-1.0) {OP};
\node[state, accepting] (no2)  at (3.6,-1.6) {OP2};
\node[state]            (nex)  at (2.2,-2.2) {$q_!$};
\node[state, accepting] (npu)  at (0,-2.2)   {PT};
\path[->]
 (q0)   edge node[sloped, above] {letra} (nid)
 (nid)  edge[loop right] node[right] {letra, dígito} ()
 (q0)   edge node[sloped, above, pos=0.6] {dígito} (nint)
 (nint) edge[loop above] node[above] {dígito} ()
 (nint) edge node[above] {.} (npt)
 (npt)  edge node[above] {dígito} (nfr)
 (nfr)  edge[loop above] node[above] {dígito} ()
 (q0)   edge node[sloped, above, pos=0.55] {aspas} (ns)
 (ns)   edge[loop above] node[above] {outro} ()
 (ns)   edge node[above] {aspas} (nsf)
 (q0)   edge node[sloped, above, pos=0.6] {< > =} (no1)
 (no1)  edge node[sloped, above] {=} (no2)
 (q0)   edge node[sloped, below, pos=0.6] {!} (nex)
 (nex)  edge node[sloped, below] {=} (no2)
 (q0)   edge node[left] {; , . \{ \} ( )} (npu);
\end{tikzpicture}
\end{adjustbox}
\caption{AFD do analisador léxico (espaços e comentários \texttt{\#} omitidos).
Em ID, o lexema é reclassificado como palavra reservada, booleano, enumerado ou
identificador; em INT e FRA, uma unidade opcional é anexada ao \textit{token}
\texttt{NUM}. Caracteres sem transição a partir de $q_0$, ou uma quebra de
linha em $q_s$, geram \textit{token} \texttt{ERRO}.}
\label{fig:afd}
\end{figure}

\subsubsection{Analisador sintático}
\label{sec:parser}
O módulo \texttt{parser} tem uma função por não terminal da
Listagem~\ref{lst:ebnf}; cada escolha de produção consulta apenas o
\textit{token} corrente, segundo os conjuntos da Tabela~\ref{tab:select}. As
repetições EBNF tornam-se laços, e a precedência
\texttt{NOT}~$>$~\texttt{AND}~$>$~\texttt{OR}, com associação à esquerda, é
dada pela própria estratificação \texttt{expr}/\texttt{conj}/\texttt{neg}. O
resultado é uma AST de nós imutáveis e tipados (\texttt{Regra},
\texttt{Comparacao}, \texttt{AcaoSet} etc.), cada um com sua posição no fonte.
Na recuperação em modo pânico, um erro descarta a declaração corrente e a
análise recomeça no próximo \texttt{POLICY}, \texttt{RULE} ou fim de arquivo,
isto é, em FIRST(\texttt{declaracao})~$\cup$~\{\texttt{\$}\}. Um
\texttt{END} ausente produz, por isso, um único erro, e a regra seguinte é
analisada normalmente.

\subsubsection{Tabela de símbolos e regras de escopo}
\label{sec:tabsimb}
O módulo \texttt{modelo} constrói a tabela de símbolos a partir de
\texttt{agroview.toml}. Cada entrada associa um nome a um domínio tipado:
booleano, enumerado, numérico (com unidade, lista de valores admissíveis e
limites mínimo e máximo) ou campo escalar; superfícies têm ainda uma
criticidade (\texttt{CRITICA}, \texttt{ALTA}, \texttt{MEDIA} ou
\texttt{BAIXA}), e campos de telemetria podem declarar um instrumento de
medição. O modelo é validado ao ser carregado: um atributo que a linguagem não
conseguiria escrever, por ser palavra reservada ou estar em caixa alta, é
rejeitado ali, e não quando uma regra tentar usá-lo. Como a DRESS-GL não
permite declarar variáveis, todos os nomes resolvem-se em um escopo global
fornecido pelo modelo, e o contexto sintático seleciona o espaço de nomes:
\texttt{view "X"} consulta as superfícies, \texttt{pass "X"} os passos,
\texttt{instances "X"} as instâncias, e uma referência sem qualificador, como
\texttt{drone.battery}, a telemetria. Nomes de regra e de política formam um
espaço de nomes próprio do programa, no qual duplicatas são rejeitadas.

\subsubsection{Verificações semânticas}
\label{sec:verif}
O módulo \texttt{semantica} percorre a AST e aplica as verificações da
Tabela~\ref{tab:verificacoes}. V2 e V3 desempenham o papel de verificação de
tipos: o tipo de cada referência é o seu domínio, e o literal comparado ou
atribuído precisa ser compatível em classe, unidade e valor. Unidades não são
convertidas: o literal deve usar exatamente a unidade do modelo, o que evita
erros de escala entre, por exemplo, milissegundos e segundos.

V6 e V7 dependem de satisfatibilidade. Como cada átomo compara uma única
variável a uma constante, o valor-verdade de um átomo só muda nas constantes
da expressão; basta testar um representante por região do domínio: o domínio
inteiro, para booleanos e enumerados; a lista de valores admissíveis, quando
existe; e, para numéricos contínuos, as constantes, os pontos médios entre
constantes consecutivas, um ponto além de cada extremo e os limites do modelo.
Avaliar a expressão em todas as combinações decide a satisfatibilidade de forma
exata. O custo cresce exponencialmente com o número de variáveis distintas,
pequeno em uma regra; acima de $10^5$ combinações, a verificação emite aviso em
vez de decidir. V6 considera apenas atributos enumerados, pois numéricos e
booleanos têm composição monotônica (Seção~\ref{sec:backend}).

\begin{table*}[t]
\caption{Verificações semânticas e casos cobertos pelos testes de semântica}
\label{tab:verificacoes}
\centering
\scriptsize
\setlength{\tabcolsep}{3pt}
\begin{tabularx}{\textwidth}{@{}l L >{\raggedright\arraybackslash}p{0.31\textwidth} c c@{}}
\toprule
\textbf{Cód.} & \textbf{Condição de rejeição} & \textbf{Exemplo rejeitado} &
\textbf{Rejeita} & \textbf{Aceita} \\
\midrule
V1 & telemetria, superfície, atributo, passo ou instância inexistente no modelo;
\texttt{SET} sobre passo & \texttt{GIVEN drone.altitude < 10 m} & 5 & -- \\
\addlinespace
V2 & unidade do literal diferente da do modelo; fator de \texttt{SCALE} com
unidade & \texttt{SET view "AGRO".hz = 2;} & 6 & -- \\
\addlinespace
V3 & valor fora dos limites, da lista admissível ou do enumerado; operador de
ordem sobre enumerado; campo escalar fora de \texttt{HIGHLIGHT}; fator fora de
$(0,1]$; \texttt{reserve} $\geq$ \texttt{budget} &
\texttt{SCALE instances "crops" BY 1.5;} & 11 & -- \\
\addlinespace
V4 & desativação de superfície \texttt{CRITICA} &
\texttt{SET view "FPV".enabled = false;} & 1 & 2 \\
\addlinespace
V5 & junção sobre passo injetado pela própria regra; \texttt{AROUND} ou
\texttt{AFTER} sobre passo suprimido; \texttt{INJECT} relativo a si mesmo ou a
passo suprimido; dois \texttt{INJECT} na mesma âncora e fase &
\texttt{WHEN AROUND pass "TerrainPass"} com \texttt{SUPPRESS pass "TerrainPass"}
& 6 & 1 \\
\addlinespace
V6 & regras de mesma prioridade, simultaneamente satisfazíveis, que atribuem
valores distintos a um atributo enumerado; dois valores na mesma regra &
duas regras escrevendo \texttt{view "MAP".mode} & 2 & 3 \\
\addlinespace
V7 & \texttt{GIVEN} insatisfazível no domínio do modelo &
\texttt{drone.battery < 10 \% AND} \texttt{drone.battery > 20 \%} & 4 & 3 \\
\addlinespace
V8 & \texttt{ASSERT} sobre telemetria sem instrumento declarado &
\texttt{ASSERT drone.battery > 5 \%} & 2 & -- \\
\midrule
C1 & \texttt{SET} que contradiz o \texttt{ASSERT} da própria regra &
\texttt{SET view "FPV".scale = 0.5} com \texttt{ASSERT view "FPV".scale >= 0.75}
& 1 & 1 \\
\addlinespace
C2 & mais de uma \texttt{POLICY}; parâmetro repetido; \texttt{POLICY} sem
\texttt{budget}; regra duplicada &
\texttt{POLICY Politica \{ cooldown = 10 frames; \}} & 4 & -- \\
\bottomrule
\end{tabularx}
\end{table*}

As verificações C1 e C2 complementam a numeração V1--V8 definida nos
requisitos.

\subsubsection{Representação intermediária e motor de execução}
\label{sec:backend}
O módulo \texttt{interpretador} só gera a RI para programas sem erros. Cada
regra torna-se uma \texttt{RegraExecutavel} com nome, prioridade, junção,
condição, contrato opcional e uma tupla de efeitos imutáveis
(\texttt{EfeitoSet}, \texttt{EfeitoScale}, \texttt{EfeitoInject},
\texttt{EfeitoSuppress}, \texttt{EfeitoHighlight}, \texttt{EfeitoMark}), cada
um com a linha de origem para rastreabilidade. A condição associa cada átomo ao
nome da leitura que o hospedeiro deve fornecer (\texttt{drone.battery},
\texttt{frame\_budget}, \texttt{view:FPV.scale}); uma leitura ausente gera erro
de execução sem alterar o estado do motor.

O módulo \texttt{motor} executa um passo por quadro: avalia os
\texttt{ASSERT}s das regras aplicadas no quadro anterior; avalia cada
\texttt{GIVEN}; ativa ou desativa regras, respeitando o \texttt{cooldown}
(uma regra que mudou de estado só pode mudar de novo após \texttt{cooldown}
quadros, e a mudança pedida antes disso é registrada como retida); e compõe os
efeitos das regras ativas em um plano, segundo a Tabela~\ref{tab:composicao}.
As regras são percorridas por prioridade decrescente, as sem prioridade por
último e os empates pela linha no fonte; como essa ordenação é feita pelo
próprio motor, o plano não depende da ordem de ativação, e compor uma regra
consigo mesma é idempotente.

\begin{table}[htbp]
\caption{Composição de efeitos entre regras ativas}
\label{tab:composicao}
\centering
\footnotesize
\begin{tabularx}{\columnwidth}{@{}>{\raggedright\arraybackslash}p{0.36\columnwidth} L@{}}
\toprule
\textbf{Efeito sobre o mesmo alvo} & \textbf{Regra de composição} \\
\midrule
\texttt{SET} numérico & mínimo \\
\texttt{SET} booleano & conjunção \\
\texttt{SET} enumerado & vence a primeira na ordem; conflito registrado \\
\texttt{SCALE} & menor fator \\
\texttt{INJECT}/\texttt{SUPPRESS} do mesmo passo & vence a primeira na ordem;
conflito registrado \\
\texttt{HIGHLIGHT}, \texttt{MARK} & união sem repetição \\
\bottomrule
\end{tabularx}
\end{table}

O motor não conhece o hospedeiro: recebe um dicionário de leituras e devolve um
registro imutável com regras ativas, ativadas, desativadas e retidas,
conflitos, violações de contrato e plano. O \texttt{ASSERT} de uma regra
aplicada no quadro $q$ é verificado com as leituras do quadro $q+1$, que medem
o efeito do plano de $q$; uma violação é registrada e não desfaz a regra.

\subsubsection{Contrato do hospedeiro}
\label{sec:contrato}
O motor conversa com qualquer hospedeiro apenas pelo contrato do módulo
\texttt{hospedeiro} (D40): no quadro $q$, o hospedeiro recebe o plano calculado
no quadro $q-1$, aplica-o, desenha e devolve as leituras de $q$, entre elas
\texttt{frame\_budget} e a configuração em vigor de cada vista. Como o motor,
o \textit{trace} e as métricas dependem só desse contrato, o mesmo programa
compilado, sem nenhuma alteração, é executado nos dois hospedeiros descritos a
seguir, o que permite comparar diretamente os resultados (Q4,
Seção~\ref{sec:avaliacao}).

\subsubsection{Hospedeiro simulado}
\label{sec:hospedeiro}
O módulo \texttt{simulador} é o hospedeiro de custo modelado. Cada cenário é um
arquivo TOML que descreve as entradas de telemetria quadro a quadro, em trechos
constantes ou em rampas lineares, com ruído gaussiano opcional de semente fixa;
o mesmo cenário produz sempre o mesmo \textit{trace}. A cada quadro, o
hospedeiro aplica o plano recebido e produz as leituras do quadro seguinte: a
configuração em vigor de cada atributo é o valor do plano ou, na sua ausência,
o padrão de maior qualidade (maior valor admissível, primeiro enumerado ou
\texttt{true}), e o custo do quadro é \emph{modelado}, não medido:
\begin{equation}
c = c_0 + \sum_{p \in P} c_p\, f_p
        + \sum_{s \in S} c_s\, e_s^{2}\, \frac{h_s}{h_s^{0}},
\label{eq:custo}
\end{equation}
em que $P$ é o conjunto de passos executados (padrão, menos os suprimidos, mais
os injetados), $f_p$ é a fração de instâncias desenhadas pelo passo $p$, $S$ é
o conjunto de superfícies habilitadas, $e_s$ é a escala linear de resolução da
superfície $s$, e $h_s$ e $h_s^{0}$ são sua frequência corrente e padrão. Os
parâmetros, calibrados manualmente, estão na seção \texttt{[custo]} de
\texttt{agroview.toml}: $c_0 = 2$\,ms; passos \texttt{TerrainPass} e
\texttt{SuitabilityPass} com 3\,ms e \texttt{CropInstancesPass} com 4\,ms;
superfícies \texttt{FPV}, \texttt{MAP} e \texttt{AGRO} com 3, 1 e 2\,ms. Modo de
visualização, destaques e marcas não entram na Eq.~(\ref{eq:custo}), e o
hospedeiro simulado não modela a ordem dos passos no grafo.

\subsubsection{Hospedeiro medido AgroView-SC}
\label{sec:agroview}
O pacote \texttt{agroview} implementa o mesmo contrato com renderização real,
em CPU, com NumPy (D42), e não sobre GPU. São três superfícies: \texttt{FPV}
($320 \times 180$), \texttt{MAP} ($128 \times 128$) e \texttt{AGRO}
($96 \times 96$), nas resoluções de \texttt{scale}~=~1,0, com 200 instâncias
base de cultura. Cada superfície é desenhada na sua própria frequência, por
decimação em relação aos 60\,Hz do quadro (D43): uma superfície a 2\,Hz é
redesenhada a cada 30 quadros. O \texttt{frame\_budget} é o custo do próprio
quadro, sem média móvel (D46), para que o \texttt{ASSERT} verificado em $q+1$
meça o efeito do plano de $q$, e não de uma janela de quadros.

O relógio tem dois modos (D48). No modo \texttt{tempo}, \texttt{frame\_budget}
é o tempo de parede do quadro, medido com \texttt{perf\_counter}. No modo
\texttt{carga}, $\mathtt{frame\_budget} = \sum_p \ell_p\, k_p$, em que a carga
$\ell_p$ é uma contagem estrutural do passo $p$ (pixels $\times$ operações por
pixel) e $k_p$ é um coeficiente por unidade de carga (em ns na Tabela~\ref{tab:calibracao}), obtido por
calibração. O modo \texttt{carga} é determinístico, o que permite fixar as
métricas por teste de regressão, e próximo do tempo medido na máquina
calibrada. A calibração (\texttt{tools/calibrar\_relogio.py}) mede 540 quadros,
depois de 60 de aquecimento, e grava em \texttt{results/calibracao.toml} a mediana de cada passo
(Tabela~\ref{tab:calibracao}). Com esses coeficientes e nenhuma regra ativa, o
quadro comum, com \texttt{TerrainPass} em \texttt{FPV} e \texttt{MAP} e as 200
instâncias de cultura, vale cerca de 3,9\,ms, e o quadro em que \texttt{AGRO}
também é redesenhado (a cada 6 quadros, a 10\,Hz), cerca de 5,1\,ms. Com a
regra ativa, o \texttt{SuitabilityPass} injetado é sobreposto a \texttt{FPV}
(D47) e a \texttt{MAP}, no modo \texttt{SUITABILITY} (D45); essas duas
sobreposições somam cerca de 9,3\,ms, e o pior quadro nas gravações chega a
14,22\,ms (Tabela~\ref{tab:coppelia}).

\begin{table}[htbp]
\caption{Coeficientes do relógio por carga}
\label{tab:calibracao}
\centering
\footnotesize
\begin{tabularx}{\columnwidth}{@{}L r r@{}}
\toprule
\textbf{Passo} & \textbf{ns/unid.\ de carga} & \textbf{Amostras} \\
\midrule
\texttt{TerrainPass} (\texttt{FPV} e \texttt{MAP}) & 50,8  & 1080 \\
\texttt{CropInstancesPass}                         & 87,8  & 540  \\
\texttt{SuitabilityPass} (\texttt{AGRO})           & 1,55  & 90   \\
\bottomrule
\end{tabularx}
\end{table}

\subsubsection{Ponte de telemetria com o CoppeliaSim}
\label{sec:ponte}
A ponte (D49--D53; a D49 substitui a D26) conecta-se ao CoppeliaSim~4.10~Edu pela
API remota ZeroMQ, em modo síncrono: a física só avança quando a ponte pede
(D49). Cada passo de física tem 50\,ms e corresponde a três quadros de
60\,Hz, nos quais vale a última leitura (D50); no teste de conexão, a latência
média foi de 7,58\,ms por passo. O simulador fornece a posição da aeronave e o
tempo simulado, sobre um campo de $5 \times 5$\,m. A entrada
\texttt{terrain.below} é o solo da zona retangular da missão
(\texttt{examples/missoes/}) sob essa posição, e \texttt{drone.battery} cai
linearmente com o tempo simulado (D51). Cada passo é gravado em JSON Lines,
junto com o alvo da missão e o solo sob o alvo, o que separa o atraso da
física do atraso do \texttt{cooldown} (D53). Foram gravados dois voos,
\texttt{pouso\_agua} (1260 quadros) e \texttt{borda\_agua} (720 quadros),
descritos na Seção~\ref{sec:avaliacao}. Os experimentos não dependem do
simulador aberto: reproduzem as gravações versionadas, quadro a quadro, em
qualquer dos dois hospedeiros (D52). Duas regravações de \texttt{borda\_agua}
com a cena versionada, reaberta antes de cada voo (\texttt{rep1} e
\texttt{rep2}), coincidiram em todos os quadros, segundo
\texttt{tools/comparar\_gravacoes.py} (D58); a gravação de referência difere
delas apenas na posição, em no máximo $5{,}0 \times 10^{-8}$\,m, e gera o
mesmo \textit{trace} (D60).

\subsubsection{Ruído de percepção}
\label{sec:percepcao}
Para examinar a política sob leituras imperfeitas, a D56 aplica à posição
percebida da aeronave, em cada eixo $x$ e $y$, um erro Gauss-Markov de
primeira ordem com desvio padrão estacionário \texttt{erro\_pos\_m} (0; 0,02;
0,05 e 0,1\,m) e constante de tempo de 0,5\,s, atualizado a 10\,Hz e mantido
entre amostras nos quadros de 60\,Hz; a física não é alterada. Como
\texttt{terrain.below} é recalculado a partir da posição percebida, o erro faz
a leitura divergir do terreno real perto da borda da água. O erro depende só
da semente e do índice do quadro; assim, uma execução com \texttt{cooldown} e
outra sem ele recebem a mesma sequência, o que isola o efeito do
\texttt{cooldown} da variação do sorteio. Os parâmetros e a hipótese foram
fixados antes da execução, e a varredura roda no hospedeiro simulado, pois as
transições não dependem do hospedeiro (Q4).

%=======================================================================
\section{Resultados e Discussão}
\label{sec:resultados-discussao}

\subsection{Compilação de um programa válido}
\label{sec:compilacao}

A Listagem~\ref{lst:ri} mostra a saída de
\texttt{python -m dressgl.interpretador examples/safe\_setdown.dgl}, que
compila a Listagem~\ref{lst:exemplo} e imprime a RI. A saída relaciona a regra
à sua junção, às leituras de que o \texttt{GIVEN} depende, a cada efeito, com a
linha de origem, e às leituras exigidas pelo \texttt{ASSERT}. Essa listagem é
a forma estática da rastreabilidade de P7; a forma dinâmica é o
\textit{trace} (Seção~\ref{sec:avaliacao}).

\begin{lstlisting}[style=saida,caption={Representação intermediária gerada para
a Listagem~\ref{lst:exemplo}.},label={lst:ri}]
RULE SafeSetdownAssist (linha 7, prioridade 90) AROUND pass "TerrainPass"
  GIVEN lê: drone.battery, terrain.below
  EfeitoInject(passo='SuitabilityPass', fase='AFTER', ancora='TerrainPass', linha=11)
  EfeitoSet(superficie='MAP', atributo='mode', valor='SUITABILITY', linha=12)
  EfeitoHighlight(extremo='ARGMAX', campo='terrain.suitability', raio_m=400, linha=13)
  EfeitoSet(superficie='AGRO', atributo='hz', valor=2, linha=14)
  EfeitoScale(recurso='instances', nome='crops', fator=0.25, linha=15)
  ASSERT lê: frame_budget, view:FPV.scale
\end{lstlisting}

\subsection{Diagnósticos de compilação}
\label{sec:diagnosticos}

A Listagem~\ref{lst:erros} contém um programa sintaticamente correto com seis
erros semânticos, e a Listagem~\ref{lst:diag}, a saída de
\texttt{python -m dressgl.semantica examples/erros\_semanticos.dgl}, sem o
prefixo com o nome do arquivo. Dois dos erros reproduzem as deficiências (i) e
(ii) da Seção~\ref{sec:relevancia}: a desativação da visão em primeira pessoa
(V4, linha~8) e a frequência sem unidade (V2, linha~9), construções que a
Listagem~\ref{lst:python} aceitaria e que só falhariam em operação. Os demais
mostram uma referência a um nome inexistente no modelo (V1), uma condição contraditória
(V7, posicionada no \texttt{AND} que a forma), uma regra cujo ponto de junção
depende de um passo que ela própria injeta (V5) e um contrato sobre uma grandeza
que o hospedeiro não mede (V8). Todos os erros são relatados em uma única
execução, ordenados por posição.

\begin{blocolst}
\begin{lstlisting}[style=dsl]
POLICY AgroViewSC {
  budget = 16.67 ms;
}

RULE PerdaGps PRIORITY 80
  GIVEN drone.altitude < 10 m
  WHEN  AROUND pass "TerrainPass"
  THEN  SET view "FPV".enabled = false;
        SET view "AGRO".hz = 2;
END

RULE PousoImpossivel PRIORITY 70
  GIVEN drone.battery < 10 % AND drone.battery > 20 %
  WHEN  AROUND pass "SuitabilityPass"
  THEN  INJECT pass "SuitabilityPass" AFTER "TerrainPass";
  ASSERT drone.battery > 5 %
END
\end{lstlisting}
\legendalst{lst:erros}{Programa com erros semânticos
(\texttt{examples/\allowbreak erros\_semanticos.dgl}).}
\end{blocolst}

\begin{lstlisting}[style=saida,caption={Diagnósticos emitidos para a
Listagem~\ref{lst:erros}.},label={lst:diag}]
6:9: erro semântico [V1]: telemetria drone.altitude não existe no modelo; disponíveis: drone.battery, frame_budget, terrain.below, terrain.suitability
8:9: erro semântico [V4]: a superfície "FPV" é CRITICA e não pode ser desativada
9:30: erro semântico [V2]: view "AGRO".hz exige unidade Hz
13:30: erro semântico [V7]: a condição GIVEN de PousoImpossivel nunca é verdadeira no domínio do modelo; a regra é inalcançável
14:9: erro semântico [V5]: WHEN AROUND pass "SuitabilityPass": o passo é injetado pela própria regra, que dependeria de si mesma para disparar
16:10: erro semântico [V8]: ASSERT usa drone.battery, que o hospedeiro não mede; declare um instrumento para ela em [telemetria]
6 erro(s), 0 aviso(s)
\end{lstlisting}

Os diagnósticos léxicos e sintáticos seguem o mesmo formato
\texttt{linha:coluna}, por exemplo
\texttt{1:6: erro léxico: caractere inesperado '@'} e
\texttt{8:9: erro sintático: esperado 'THEN', encontrado 'SUPPRESS'}.

\subsection{Testes automatizados}
\label{sec:testes}

A Tabela~\ref{tab:testes} resume os 223 testes. O comando
\texttt{python -m pytest -q} executa os 222 rápidos, em cerca de 10\,s; o teste
marcado como \texttt{lento}, que regenera \texttt{results/metricas\_coppelia.csv}
e o compara com a versão versionada em cerca de 1\,min, é executado por
\texttt{python -m pytest -q -m lento} (D59). Os testes de \texttt{test\_artigo}
fixam as Listagens~\ref{lst:ri} e~\ref{lst:diag}: se o compilador mudar, o
texto do artigo precisa mudar junto. Os de
\texttt{test\_experimentos} fixam os números da Tabela~\ref{tab:metricas}, e
os de \texttt{test\_regressao\_csv}, os de
\texttt{results/metricas\_coppelia.csv} e
\texttt{results/varredura\_ruido.csv}, que alimentam as
Tabelas~\ref{tab:coppelia} e~\ref{tab:varredura}.

\begin{table}[htbp]
\caption{Testes automatizados por grupo}
\label{tab:testes}
\centering
\footnotesize
\setlength{\tabcolsep}{4pt}
\begin{tabularx}{\columnwidth}{@{}L r@{}}
\toprule
\textbf{Grupo (arquivos)} & \textbf{Testes} \\
\midrule
Léxico, gramática e sintático (\texttt{test\_lexer}, \texttt{test\_gramatica},
\texttt{test\_parser}) & 36 \\
Semântica e modelo (\texttt{test\_semantica}) & 57 \\
RI e motor (\texttt{test\_interpretador}, \texttt{test\_motor}) & 25 \\
Hospedeiro simulado e Tabela~\ref{tab:metricas} (\texttt{test\_simulador},
\texttt{test\_experimentos}) & 29 \\
Contrato e núcleo sem NumPy (\texttt{test\_hospedeiro},
\texttt{test\_nucleo\_sem\_numpy}) & 5 \\
AgroView-SC: hospedeiro, ordem dos passos, render e relógio
(\texttt{test\_agroview\_hospedeiro}, \texttt{test\_agroview\_ordem},
\texttt{test\_agroview\_render}, \texttt{test\_relogio\_d48}) & 40 \\
Ponte e experimentos com o CoppeliaSim (\texttt{test\_ponte\_coppelia},
\texttt{test\_experimentos\_coppelia}) & 14 \\
Ruído de percepção (\texttt{test\_percepcao}) & 12 \\
Regressão dos CSVs (\texttt{test\_regressao\_csv}, inclui 1 \texttt{lento}) & 2 \\
Artigo e pacote (\texttt{test\_artigo}, \texttt{test\_pacote}) & 3 \\
\midrule
\textbf{Total} & \textbf{223} \\
\bottomrule
\end{tabularx}
\end{table}

\subsection{Avaliação experimental}
\label{sec:avaliacao}

A avaliação examina o comportamento, em execução, das políticas compiladas,
orientada por cinco questões:
\begin{itemize}
\item[\textbf{Q1}] As compensações declaradas em \texttt{THEN} mantêm o custo
do quadro dentro do orçamento, em comparação com uma política que apenas reage
ao evento sem compensá-lo? (F5, F6, P6)
\item[\textbf{Q2}] Qual é o efeito do \texttt{cooldown} sobre a oscilação de
ativação e sobre o atraso de resposta? (F1, P6)
\item[\textbf{Q3}] Com mais de uma regra ativa, a composição produz o plano
previsto pela Tabela~\ref{tab:composicao} e registra os conflitos de forma
rastreável? (P6, P7)
\item[\textbf{Q4}] As conclusões de Q1 dependem do hospedeiro, isto é, se
mantêm quando o custo é medido no AgroView-SC em vez de modelado? (P6, P8)
\item[\textbf{Q5}] Com telemetria de origem física e ruído de percepção, o
\texttt{cooldown} continua reduzindo as transições, e a que custo? (F1, P6, P8)
\end{itemize}

Q1--Q3 usam o hospedeiro simulado e os cenários sintéticos; Q4 usa as duas
gravações do CoppeliaSim nos dois hospedeiros, com o AgroView-SC no relógio
\texttt{carga}; Q5 usa a gravação \texttt{borda\_agua}, sem e com erro de
posição. A Tabela~\ref{tab:condicoes}
relaciona as sete condições em que o cenário de bateria em reserva sobre
superfície imprópria é examinado e a evidência de cada uma.

\begin{table}[htbp]
\caption{Condições de avaliação do cenário único e evidência correspondente}
\label{tab:condicoes}
\centering
\scriptsize
\setlength{\tabcolsep}{3pt}
\begin{tabularx}{\columnwidth}{@{}>{\raggedright\arraybackslash}p{0.25\columnwidth} L c@{}}
\toprule
\textbf{Condição} & \textbf{Evidência} & \textbf{Q} \\
\midrule
Operação normal & quadros anteriores à ativação em \texttt{listagem3}
(0--219) e nas duas gravações & Q1, Q4 \\
\addlinespace
Ativação e compensação & DSL $\times$ \textit{baseline} em \texttt{listagem3}
e nas duas gravações, nos dois hospedeiros & Q1, Q4 \\
\addlinespace
Oscilação & \texttt{oscilacao} (sintético) e \texttt{borda\_agua} (física) &
Q2, Q5 \\
\addlinespace
Ruído no custo & \texttt{listagem3\_ruido} & Q1 \\
\addlinespace
Ruído na percepção & erro de posição da D56 em \texttt{borda\_agua}: 4 níveis (3 não nulos)
$\times$ 10 sementes, 40 pares & Q5 \\
\addlinespace
Composição de regras & múltiplas regras nos três cenários simulados e nas duas
gravações & Q3, Q4 \\
\addlinespace
Rejeição estática & Listagens~\ref{lst:erros} e~\ref{lst:diag} (V1, V2, V4,
V5, V7, V8) & -- \\
\bottomrule
\end{tabularx}
\end{table}

\textbf{Programas comparados.}
(a)~\textit{DSL}: a Listagem~\ref{lst:exemplo};
(b)~\textit{baseline}: a mesma regra, com a mesma \texttt{POLICY} e as mesmas
cláusulas \texttt{GIVEN}, \texttt{WHEN} e \texttt{ASSERT}, mas sem as
compensações (\texttt{SET view "AGRO".hz} e \texttt{SCALE}), o que isola o
efeito destas;
(c)~\textit{DSL sem cooldown}: o programa (a) com \texttt{cooldown} igual a
zero, como ablação;
(d)~\textit{múltiplas regras}: o programa (a) acrescido de uma regra sintética,
\texttt{EconomiaBateria} (\texttt{PRIORITY 50}), ativa com bateria abaixo de
25\,\%, que escreve \texttt{view "MAP".mode = NOMINAL},
\texttt{view "AGRO".hz = 1 Hz} e \texttt{SCALE instances "crops" BY 0.5},
construída para provocar escritas concorrentes.

\textbf{Cenários sintéticos.}
\texttt{listagem3}: 300 quadros; a bateria decresce linearmente de 40\,\% para
10\,\%, e a aeronave sobrevoa solo exposto (\texttt{BARE\_SOIL}), exceto entre
os quadros 180 e 259, em que sobrevoa água. A bateria cruza 18\,\% no quadro
220; o \texttt{GIVEN} da Listagem~\ref{lst:exemplo} é verdadeiro nos quadros
220 a 259 (40 quadros, 13,33\,\% do total);
\texttt{listagem3\_ruido}: o mesmo cenário, com ruído gaussiano aditivo sobre o
custo, sorteado a cada quadro com $\sigma = 0{,}5$\,ms e semente 42; como a
semente é fixa, a mesma sequência de ruído é aplicada a todos os programas;
\texttt{oscilacao}: 150 quadros sobre água, com a bateria em 19\,\% e quedas a
17\,\% nos quadros 50--54, 60--64 e de 70 em diante.

\textbf{Cenários físicos (gravações do CoppeliaSim).}
\texttt{pouso\_agua}: 1260 quadros (21\,s a 60\,Hz), em que a condição do
\texttt{GIVEN} se torna verdadeira uma vez, com uma ativação e uma desativação
da regra; \texttt{borda\_agua}: 720 quadros (12\,s), em que a aeronave cruza a
borda de um corpo d'água e \texttt{terrain.below} muda seis vezes, o que produz
três ativações e três desativações. Nas duas gravações, as mudanças vêm da
dinâmica do voo na cena, e não de trechos programados como em
\texttt{oscilacao}.

\textbf{Métricas.} Por par programa--cenário: porcentagem de quadros acima do
orçamento de 16,67\,ms; pior quadro; ativações e desativações; retenções pelo
\texttt{cooldown}; violações de \texttt{ASSERT}; conflitos de composição; e
atraso de ativação, em quadros, entre o quadro em que a condição se torna
verdadeira e o quadro em que a regra é ativada.

\textbf{Reprodutibilidade.} A Tabela~\ref{tab:metricas} é gerada por
\texttt{python -m dressgl.experimentos}, que executa a matriz de $4 \times 3$
combinações no hospedeiro simulado e grava \texttt{results/metricas.csv} e
\texttt{results/metricas.md}; seus números não mudaram desde o \textit{commit}
\texttt{ddd62d0} da versão~1.0. A Tabela~\ref{tab:coppelia} reproduz
\texttt{results/metricas\_coppelia.csv}, gerado por
\texttt{python -m ponte\_coppelia.experimentos} (D54), com os quatro programas sobre as duas
gravações nos dois hospedeiros, e a Tabela~\ref{tab:varredura} resume
\texttt{results/varredura\_ruido.csv}, gerado por
\texttt{python tools/varredura\_ruido.py} sobre
\texttt{examples/missoes/margem\_ruido.toml} (4 níveis de erro $\times$ 10
sementes $\times$ 2 programas, 80 execuções). Os três CSVs são versionados e
fixados por testes de regressão (Seção~\ref{sec:testes}); os \textit{traces}
JSON Lines não são versionados e são regeneráveis pelos mesmos comandos. Todos
os resultados correspondem à \textit{tag} \texttt{v3.0}.

\begin{table*}[t]
\caption{Métricas por programa e cenário no hospedeiro simulado, com custo
modelado pela Eq.~(\ref{eq:custo})}
\label{tab:metricas}
\centering
\footnotesize
\setlength{\tabcolsep}{5pt}
\begin{adjustbox}{max width=\textwidth}
\begin{tabular}{@{}llrrcrrrc@{}}
\toprule
\textbf{Programa} & \textbf{Cenário} &
\textbf{\cab{\%\\acima}} & \textbf{\cab{Pior\\(ms)}} &
\textbf{\cab{Ativ./\\desat.}} & \textbf{Retenções} &
\textbf{\cab{Viol.\\ASSERT}} & \textbf{Conflitos} &
\textbf{\cab{Atraso\\méd./máx.}} \\
\midrule
DSL              & \texttt{listagem3}        & 0,00  & 15,00 & 1/1 & 0  & 0  & 0  & 0,00/0 \\
DSL              & \texttt{listagem3\_ruido} & 0,00  & 16,32 & 1/1 & 0  & 0  & 0  & 0,00/0 \\
DSL              & \texttt{oscilacao}        & 0,00  & 15,00 & 2/1 & 10 & 0  & 0  & 2,50/5 \\
\addlinespace
Baseline         & \texttt{listagem3}        & 13,33 & 18,00 & 1/1 & 0  & 40 & 0  & 0,00/0 \\
Baseline         & \texttt{listagem3\_ruido} & 13,33 & 18,79 & 1/1 & 0  & 40 & 0  & 0,00/0 \\
Baseline         & \texttt{oscilacao}        & 59,33 & 18,00 & 2/1 & 10 & 89 & 0  & 2,50/5 \\
\addlinespace
DSL s/ cooldown  & \texttt{listagem3}        & 0,00  & 15,00 & 1/1 & 0  & 0  & 0  & 0,00/0 \\
DSL s/ cooldown  & \texttt{listagem3\_ruido} & 0,00  & 16,32 & 1/1 & 0  & 0  & 0  & 0,00/0 \\
DSL s/ cooldown  & \texttt{oscilacao}        & 0,00  & 15,00 & 3/2 & 0  & 0  & 0  & 0,00/0 \\
\addlinespace
Múltiplas regras & \texttt{listagem3}        & 0,00  & 15,00 & 2/1 & 0  & 0  & 40 & 0,00/0 \\
Múltiplas regras & \texttt{listagem3\_ruido} & 0,00  & 16,32 & 2/1 & 0  & 0  & 40 & 0,00/0 \\
Múltiplas regras & \texttt{oscilacao}        & 0,00  & 15,00 & 3/1 & 10 & 0  & 90 & 1,67/5 \\
\bottomrule
\end{tabular}
\end{adjustbox}
\end{table*}

\begin{table}[htbp]
\caption{Custo modelado do quadro por configuração ativa}
\label{tab:custos}
\centering
\footnotesize
\begin{tabularx}{\columnwidth}{@{}L r@{}}
\toprule
\textbf{Configuração ativa} & \textbf{Custo (ms)} \\
\midrule
Nenhuma regra (configuração padrão)             & 15,0 \\
\texttt{SafeSetdownAssist} (programa DSL)       & 13,4 \\
Regra sem compensações (\textit{baseline})      & 18,0 \\
\texttt{EconomiaBateria}                        & 11,2 \\
\texttt{SafeSetdownAssist} e \texttt{EconomiaBateria} & 13,2 \\
\bottomrule
\end{tabularx}
\end{table}

\textbf{Q1: compensações e orçamento.} No cenário \texttt{listagem3}, a
política em DRESS-GL não excede o orçamento em nenhum quadro; seu pior quadro
(15,00\,ms) é o da configuração padrão, anterior à ativação. O
\textit{baseline} excede o orçamento em 13,33\,\% dos quadros: são os 40
quadros que medem o efeito da regra, com 18,00\,ms cada, e cada um registra uma
violação do contrato \texttt{frame\_budget <= 16.67 ms}. A
Tabela~\ref{tab:custos} explica a diferença. Das ações mantidas no
\textit{baseline}, apenas a injeção do \texttt{SuitabilityPass} tem custo no
modelo, de 3,0\,ms, que eleva o quadro de 15,0 para 18,0\,ms; o modo
\texttt{SUITABILITY} e o destaque não entram na Eq.~(\ref{eq:custo}). As
compensações da Listagem~\ref{lst:exemplo} retiram 4,6\,ms (3,0\,ms pela
densidade de \texttt{crops} a 0,25 e 1,6\,ms pela camada agronômica a 2\,Hz),
de modo que o quadro com a regra ativa custa 13,4\,ms, abaixo da configuração
padrão. No cenário \texttt{oscilacao}, a diferença se amplia (0\,\% contra
59,33\,\%, com 89 violações), porque a regra permanece ativa por mais tempo. O
ruído não altera as porcentagens em nenhum programa; eleva apenas o pior
quadro, para 16,32\,ms na DSL, ainda abaixo do orçamento, e para 18,79\,ms no
\textit{baseline}. Esse resultado mostra que as compensações declaradas têm o
efeito pretendido \emph{dentro do modelo} de custo, com parâmetros escolhidos
pelos autores; Q4 examina o mesmo par de programas no AgroView-SC.

\textbf{Q2: efeito do cooldown.} Em \texttt{listagem3}, as execuções com e sem
\texttt{cooldown} são idênticas: a ativação (quadro 220) e a desativação
(quadro 260) distam 40 quadros, mais que o período de 10. Em
\texttt{oscilacao}, sem \texttt{cooldown} a regra segue cada queda: ativa em 50,
60 e 70 e desativa em 55 e 65 (3 ativações e 2 desativações). Com
\texttt{cooldown} de 10 quadros, a desativação pedida em 55 fica retida até a
condição voltar a valer em 60, a regra desativa em 65 e a reativação pedida em
70 só ocorre em 75: 2 ativações, 1 desativação e 10 retenções. A contrapartida
é o atraso de ativação, que passa de zero para 2,50 quadros em média e 5 no
máximo. Como nenhuma das variantes excede o orçamento nesse cenário, o
\texttt{cooldown} atua sobre a estabilidade da configuração exibida, e não sobre
o custo; a troca entre estabilidade e tempo de resposta fica explícita e
ajustável na \texttt{POLICY}.

\textbf{Q3: composição de múltiplas regras.} Em \texttt{listagem3},
\texttt{EconomiaBateria} permanece ativa do quadro 150 até o fim, e
\texttt{SafeSetdownAssist}, entre os quadros 220 e 259. O \textit{trace}
registra um conflito por quadro de coatividade (40 em \texttt{listagem3} e 90
em \texttt{oscilacao}), todos sobre \texttt{view "MAP".mode}, atributo
enumerado sem composição monotônica: a escrita é resolvida pela prioridade
(90 contra 50), e cada registro identifica a regra vencedora e a perdedora
(Listagem~\ref{lst:trace}, extraída de
\texttt{trace\_multiplas\_\_oscilacao.jsonl}). Para os atributos numéricos, a
composição pelo mínimo dispensa resolução de conflito: o plano composto usa
\texttt{view "AGRO".hz} igual a 1\,Hz, valor de \texttt{EconomiaBateria}, e
\texttt{instances "crops"} igual a 0,25, valor de \texttt{SafeSetdownAssist},
com custo de 13,2\,ms. Após a desativação de \texttt{SafeSetdownAssist}, o
plano volta ao de \texttt{EconomiaBateria} (11,2\,ms). O orçamento é respeitado
em todos os quadros dos três cenários.

\begin{blocolst}
\begin{lstlisting}[style=saida]
{"quadro": 148,
 "ativas": ["SafeSetdownAssist", "EconomiaBateria"],
 "conflitos": [{"alvo": "MAP.mode", "vencedora": "SafeSetdownAssist",
                "perdedora": "EconomiaBateria"}],
 "violacoes": [],
 "plano": {"sets": {"MAP.mode": "SUITABILITY", "AGRO.hz": 1},
           "escalas": {"instances.crops": 0.25},
           "injecoes": [{"passo": "SuitabilityPass", "fase": "AFTER",
                         "ancora": "TerrainPass", "linha": 11}], ...},
 "leituras": {"drone.battery": 17.0, "frame_budget": 13.2,
              "view:MAP.mode": "SUITABILITY", "view:AGRO.hz": 1, ...}}
\end{lstlisting}
\legendalst{lst:trace}{Trecho do registro do quadro 148 do \textit{trace} de
múltiplas regras no cenário \texttt{oscilacao} (campos omitidos indicados por
\texttt{...}).}
\end{blocolst}

\begin{table*}[t]
\caption{Programas nas gravações do CoppeliaSim, nos dois hospedeiros
(AgroView-SC com relógio por carga)}
\label{tab:coppelia}
\centering
\footnotesize
\setlength{\tabcolsep}{5pt}
\begin{adjustbox}{max width=\textwidth}
\begin{tabular}{@{}lllrrcrr@{}}
\toprule
\textbf{Gravação} & \textbf{Hospedeiro} & \textbf{Programa} &
\textbf{\cab{\%\\acima}} & \textbf{\cab{Pior\\(ms)}} &
\textbf{\cab{Ativ./\\desat.}} & \textbf{\cab{Viol.\\ASSERT}} &
\textbf{Conflitos} \\
\midrule
\texttt{borda\_agua} & Simulado    & DSL              & 0,00  & 15,00 & 3/3 & 0   & 0   \\
\texttt{borda\_agua} & Simulado    & DSL s/ cooldown  & 0,00  & 15,00 & 3/3 & 0   & 0   \\
\texttt{borda\_agua} & Simulado    & Baseline         & 39,17 & 18,00 & 3/3 & 282 & 0   \\
\texttt{borda\_agua} & Simulado    & Múltiplas regras & 0,00  & 15,00 & 4/3 & 0   & 282 \\
\addlinespace
\texttt{borda\_agua} & AgroView-SC & DSL              & 0,00  & 14,22 & 3/3 & 0   & 0   \\
\texttt{borda\_agua} & AgroView-SC & DSL s/ cooldown  & 0,00  & 14,22 & 3/3 & 0   & 0   \\
\texttt{borda\_agua} & AgroView-SC & Baseline         & 0,00  & 14,33 & 3/3 & 0   & 0   \\
\texttt{borda\_agua} & AgroView-SC & Múltiplas regras & 0,00  & 14,22 & 4/3 & 0   & 282 \\
\addlinespace
\texttt{pouso\_agua} & Simulado    & DSL              & 0,00  & 15,00 & 1/1 & 0   & 0   \\
\texttt{pouso\_agua} & Simulado    & DSL s/ cooldown  & 0,00  & 15,00 & 1/1 & 0   & 0   \\
\texttt{pouso\_agua} & Simulado    & Baseline         & 20,48 & 18,00 & 1/1 & 258 & 0   \\
\texttt{pouso\_agua} & Simulado    & Múltiplas regras & 0,00  & 15,00 & 2/1 & 0   & 258 \\
\addlinespace
\texttt{pouso\_agua} & AgroView-SC & DSL              & 0,00  & 14,22 & 1/1 & 0   & 0   \\
\texttt{pouso\_agua} & AgroView-SC & DSL s/ cooldown  & 0,00  & 14,22 & 1/1 & 0   & 0   \\
\texttt{pouso\_agua} & AgroView-SC & Baseline         & 0,00  & 14,33 & 1/1 & 0   & 0   \\
\texttt{pouso\_agua} & AgroView-SC & Múltiplas regras & 0,00  & 14,22 & 2/1 & 0   & 258 \\
\bottomrule
\end{tabular}
\end{adjustbox}
\end{table*}

\textbf{Q4: dependência do hospedeiro.} A Tabela~\ref{tab:coppelia} executa os
quatro programas sobre as duas gravações, nos dois hospedeiros; em todas as
linhas, as retenções e o atraso de ativação são nulos. No hospedeiro simulado,
as gravações confirmam Q1: a DSL não excede o orçamento, e o \textit{baseline}
o excede em 39,17\,\% dos quadros de \texttt{borda\_agua} e em 20,48\,\% dos de
\texttt{pouso\_agua}, com uma violação de \texttt{ASSERT} em cada um desses
quadros (282 e 258). No AgroView-SC, nenhum programa excede o orçamento ou viola
contrato: o pior quadro é de 14,22\,ms na DSL e de 14,33\,ms no
\textit{baseline}. A direção do efeito se mantém, já que o \textit{baseline}
continua mais caro, mas a diferença no pior quadro cai de 3,00\,ms para
0,11\,ms, e as compensações deixam de ser necessárias para cumprir o orçamento.
A razão está no nível de custo, e não no custo do passo injetado: no
AgroView-SC, o quadro sem regra ativa custa cerca de 3,9\,ms, contra 15,0\,ms
no modelo, e a injeção, sobreposta a \texttt{FPV} e \texttt{MAP}, acrescenta
cerca de 9,3\,ms, de modo que mesmo o \textit{baseline} fica abaixo do
orçamento. As compensações quase não alteram o pior quadro: o
\texttt{CropInstancesPass} custa cerca de 0,16\,ms neste renderizador, contra
4\,ms no modelo, e reduzi-lo a 0,25 retira os 0,11\,ms observados; a camada
\texttt{AGRO} a 2\,Hz é redesenhada com menos frequência, mas, por decimação
(D43), cada redesenho custa o mesmo, e o pico não muda. Os 9,3\,ms das
sobreposições usam o coeficiente do \texttt{SuitabilityPass} medido apenas em
\texttt{AGRO}: um erro de cerca de 25\,\% nesse coeficiente consumiria a
margem de 2,34\,ms do \textit{baseline} e a de 2,45\,ms da DSL, estimadas,
além disso, pelo relógio \texttt{carga}, que tende a ficar um pouco abaixo do
tempo de parede (Seção~\ref{sec:ameacas}). O que se transfere entre hospedeiros é, assim, o
comportamento do motor, e não a conclusão sobre o orçamento: ativações,
desativações e conflitos coincidem nos dois hospedeiros (por exemplo, 4/3 e 282
conflitos no programa de múltiplas regras em \texttt{borda\_agua}), porque
dependem só da telemetria. A conclusão de Q1 vale dentro do modelo de custo, e
não para este renderizador.

\begin{table}[htbp]
\caption{Varredura do erro de posição (D56) na gravação borda\_agua: médias
de 10 sementes por nível}
\label{tab:varredura}
\centering
\footnotesize
\setlength{\tabcolsep}{3pt}
\begin{adjustbox}{max width=\columnwidth}
\begin{tabular}{@{}lrrrrc@{}}
\toprule
\textbf{\cab{Erro\\(m)}} & \textbf{\cab{Quadros\\diverg.}} &
\textbf{\cab{Trocas e trans.\\s/ cool.}} &
\textbf{\cab{Trans.\\c/ cool.}} & \textbf{\cab{Reten-\\ções}} &
\textbf{\cab{Pares c/\\redução}} \\
\midrule
0,00 & 0,0   & 6,0  & 6,0  & 0,0  & 0/10 \\
0,02 & 102,9 & 13,2 & 12,4 & 12,7 & 3/10 \\
0,05 & 160,8 & 16,8 & 14,4 & 16,9 & 6/10 \\
0,10 & 193,5 & 18,9 & 16,7 & 26,0 & 7/10 \\
\bottomrule
\end{tabular}
\end{adjustbox}
\end{table}

\textbf{Q5: telemetria física e cooldown sob ruído.} Sem erro de posição, as
seis trocas de terreno de \texttt{borda\_agua} produzem 3 ativações e 3
desativações com e sem \texttt{cooldown}, sem retenções
(Tabela~\ref{tab:coppelia}): as trocas físicas distam mais que 10 quadros, e o
\texttt{cooldown} não atua. A oscilação rápida surge com o erro de posição da
D56 (Tabela~\ref{tab:varredura}). Perto da borda, a leitura percebida diverge
da real em 14,3\,\% a 26,9\,\% dos 720 quadros, e as trocas percebidas sobem
de 6 para 13,2, 16,8 e 18,9, em média. Sem \texttt{cooldown}, a regra segue
cada troca percebida, e as transições (ativações mais desativações) igualam as
trocas. Com \texttt{cooldown} de 10 quadros, as transições caem 6,1\,\%,
14,3\,\% e 11,6\,\% nos três níveis; descontadas as seis transições reais, a
queda das espúrias é de 11\,\%, 22\,\% e 17\,\%. Dos 30 pares com erro, o
\texttt{cooldown} reduziu as transições em 16 e as igualou em 14, sem
aumentá-las em nenhum. Pelo teste do sinal unilateral sobre as ativações pareadas, a diferença é
significativa a 5\,\% nos níveis de 0,05 e 0,1\,m ($p = 0{,}016$ e
$0{,}008$), mas não no de 0,02\,m ($p = 0{,}125$). O efeito é parcial: mesmo com \texttt{cooldown}, restam
de 6,4 a 10,7 transições espúrias por execução. O custo aparece nas retenções,
12,7, 16,9 e 26,0 mudanças adiadas por execução, que podem incluir trocas
reais; a varredura não mede o atraso resultante. Nenhuma das 80 execuções
violou contrato ou registrou conflito. A troca entre estabilidade e tempo de
resposta de Q2 permanece com percepção imperfeita, mas um \texttt{cooldown} de
10 quadros, sozinho, não elimina as transições espúrias.

\subsection{Verificação dos requisitos}
\label{sec:requisitos-verif}

A Tabela~\ref{tab:requisitos} confronta os requisitos da
Seção~\ref{sec:requisitos} com a evidência obtida. Quanto à pergunta de se a
DSL simplifica a tarefa no domínio, a evidência deste trabalho é a capacidade
de verificação: erros que a formulação da Listagem~\ref{lst:python} aceitaria,
como os da Listagem~\ref{lst:erros}, são rejeitados antes da execução, e a
interação entre regras é decidida por regras de composição explícitas, e não
pela ordem de registro. A execução do mesmo programa compilado em dois
hospedeiros mostra, além disso, que a política é independente da
implementação do compositor. A simplificação para o especialista, no sentido
de legibilidade e facilidade de modificação, não foi medida
(Seção~\ref{sec:ameacas}).

\begin{table*}[t]
\caption{Verificação dos requisitos}
\label{tab:requisitos}
\centering
\scriptsize
\setlength{\tabcolsep}{3pt}
\begin{tabularx}{\textwidth}{@{}l >{\raggedright\arraybackslash}p{0.14\textwidth} L@{}}
\toprule
\textbf{Req.} & \textbf{Situação} & \textbf{Evidência} \\
\midrule
F1 & atendido & \texttt{POLICY} com \texttt{budget}, \texttt{reserve} e \texttt{cooldown} verificados (V2, V3, C2) e usados pelo motor; Q2 e Q5 \\
F2 & atendido & gabarito completo na gramática (Listagem~\ref{lst:ebnf}) e na AST; prioridade inteira verificada no sintático \\
F3 & atendido & \texttt{AND}, \texttt{OR}, \texttt{NOT}, \texttt{IN} e unidades; testes de precedência e de V1--V3 em \texttt{GIVEN} \\
F4 & parcial & junção analisada, verificada (V5) e preservada na RI; em nenhum dos dois hospedeiros o ponto de junção altera a execução \\
F5 & atendido, com ressalva & as seis ações compilam e são compostas; \texttt{SET}, \texttt{SCALE} e \texttt{INJECT} alteram o custo nos dois hospedeiros; \texttt{HIGHLIGHT} e \texttt{MARK} chegam ao plano, mas não são desenhados \\
F6 & atendido & \texttt{ASSERT} verificado estaticamente (V8, C1) e monitorado em execução no quadro $q+1$; 40 e 89 violações no \textit{baseline} simulado, 282 e 258 nas gravações; nenhuma no AgroView-SC (Q4) \\
F7 & atendido & a gramática não tem construções de laço, recursão, função ou aritmética \\
\addlinespace
P1 & atendido & AFD da Fig.~\ref{fig:afd}; testes do analisador léxico \\
P2 & atendido & descendente recursivo; propriedade LL(1) verificada por teste; recuperação em modo pânico \\
P3 & atendido & tabela de símbolos derivada de \texttt{agroview.toml}, validada na carga; resolução por espaço de nomes \\
P4 & atendido & V1--V8 com posição (Tabela~\ref{tab:verificacoes}, Listagem~\ref{lst:diag}) \\
P5 & atendido & RI de condição e efeitos imutáveis (Listagem~\ref{lst:ri}); só é gerada para programas sem erros \\
P6 & atendido nos dois hospedeiros & motor com composição, \texttt{cooldown} e \texttt{ASSERT}; Q1--Q5; não integrado a um compositor sobre GPU \\
P7 & atendido & RI com linha de origem de cada efeito; \textit{trace} com regra vencedora e perdedora de cada conflito (Listagem~\ref{lst:trace}) \\
P8 & atendido & contrato D40 com duas implementações; núcleo sem NumPy verificado em subprocesso; ponte com gravações de reprodução determinística (D58); Q4 e Q5 \\
\bottomrule
\end{tabularx}
\end{table*}

\subsection{Ameaças à validade}
\label{sec:ameacas}

\textbf{Validade de construção.} No hospedeiro simulado, os custos são
modelados e calibrados manualmente. No AgroView-SC, são medidos, mas sobre um
renderizador em CPU e em resolução reduzida (D42), e não sobre um compositor de
produção em GPU; o custo relativo dos passos difere, por isso, do que se
esperaria em GPU. O relógio por carga é linear e não tem termo fixo por passo,
de modo que passos com custo fixo relevante, como o
\texttt{CropInstancesPass}, ficam subestimados quando a carga cai; ele mede os
passos, e não o intervalo entre eles, e tende a ficar um pouco abaixo do
relógio \texttt{tempo} (D48). O hospedeiro simulado usa o custo médio das
superfícies decimadas, e o AgroView-SC, o custo de cada quadro, com os picos
dos quadros em que elas são redesenhadas (D43, D46); por isso, o pior quadro
só é comparável entre programas no mesmo hospedeiro. O \texttt{SuitabilityPass} foi calibrado apenas
na superfície \texttt{AGRO}; as variantes desenhadas sobre \texttt{MAP} e
\texttt{FPV} usam o mesmo coeficiente sem medição própria e respondem por cerca
de 9,3 dos 14,22\,ms do pior quadro da DSL em Q4; a conclusão de que nenhum
programa excede o orçamento no AgroView-SC depende, portanto, desse
coeficiente. Os resultados do
AgroView-SC usam o relógio \texttt{carga} e não foram repetidos no relógio
\texttt{tempo}. O \texttt{frame\_budget} sem média (D46) torna o
\texttt{ASSERT} sensível a quadros isolados. O \textit{baseline} foi construído
pelos autores a partir da própria política em DRESS-GL; não é uma
implementação imperativa independente, e a Listagem~\ref{lst:python} não foi
executada. A independência da ordem de composição foi testada para pares de
regras nas duas ordens, mas não de forma exaustiva. V4 cobre apenas a
desativação de superfícies críticas: uma redução de resolução da visão em
primeira pessoa por \texttt{SCALE resolution} é aceita pelo compilador e só é
detectada em execução, se houver \texttt{ASSERT} correspondente.

\textbf{Validade interna.} A regra \texttt{EconomiaBateria} é sintética, criada
para exercitar a composição, e não provém de procedimento operacional real. Os
cenários produzem poucas ativações. O efeito do \texttt{cooldown} foi observado
em um cenário oscilante sintético e, com telemetria física, apenas sob erro de
posição: sem erro, as trocas de terreno de \texttt{borda\_agua} distam mais que
10 quadros e não o acionam. A varredura usa uma única gravação, quatro níveis (três não nulos) de
erro e dez sementes por nível, e não mede o atraso causado pelas retenções.
Pontos de junção, destaques e marcas não influenciam nenhum dos hospedeiros. O
atraso físico (D53), contado do primeiro quadro com o alvo da missão sobre a água ao primeiro quadro com a aeronave sobre a água e registrado em \texttt{results/metricas\_coppelia.csv} (36 quadros
em \texttt{borda\_agua} e 24 em \texttt{pouso\_agua}), é o mesmo para todos os
programas e para os dois hospedeiros; depende, portanto, só da gravação, e não
da política. A gravação de referência de \texttt{borda\_agua} difere das regravações apenas
na posição, em no máximo $5{,}0 \times 10^{-8}$\,m, e gera o mesmo
\textit{trace} (D60); a cena de \texttt{pouso\_agua} não está versionada, e
essa gravação não pode ser refeita. Como a geração de uma gravação depende do
simulador, a evidência reproduzível é a gravação versionada, e não a sua
regravação.

\textbf{Validade externa.} A telemetria física vem de uma cena simplificada de
$5 \times 5$\,m no CoppeliaSim, reproduzida a partir de gravações, sem laço
fechado entre a política e o voo e sem aeronave real. A bateria cai
linearmente com o tempo simulado, sem relação com o empuxo (D51); a
classificação do solo é exata, dada a posição; e o ruído de percepção não foi
calibrado contra um receptor real (D56). Apenas um dos seis
cenários excepcionais da Seção~\ref{sec:requisitos} foi avaliado, por decisão
de escopo (D57). As medições valem para a máquina de
\texttt{results/ambiente.json}. Não foi conduzido o estudo de legibilidade com
participantes, de modo que a hipótese de que a notação é mais auditável pelo
especialista permanece não verificada~\cite{rodrigues2017}.

%=======================================================================
\section{Conclusão}
\label{sec:conclusao}

Este trabalho apresentou a DRESS-GL, linguagem específica de domínio externa e
não Turing-completa para políticas de degradação gráfica em estações de
controle de VANT, e seu processador, que cobre as análises léxica, sintática e
semântica, compila as regras em uma representação intermediária de condição e
efeitos e as executa por quadro. A gramática teve a propriedade LL(1)
verificada automaticamente; as oito verificações semânticas rejeitam, antes da
execução, construções que uma linguagem de propósito geral aceitaria, como a
desativação de uma superfície crítica ou uma frequência sem unidade. Um
contrato de hospedeiro permitiu executar o mesmo programa compilado sobre um
hospedeiro de custo modelado e sobre o AgroView-SC, de custo medido, alimentado
por telemetria gravada no CoppeliaSim. No hospedeiro simulado, as compensações
declaradas mantiveram o quadro no orçamento, o \texttt{cooldown} tornou
explícita a troca entre estabilidade e tempo de resposta, e a composição de
regras registrou cada conflito de forma rastreável. No AgroView-SC, com o relógio por carga, nenhum
programa excedeu o orçamento, nem mesmo o \textit{baseline}: a vantagem das
compensações depende do modelo de custo, e a margem restante está condicionada a um
coeficiente do \texttt{SuitabilityPass} medido só em \texttt{AGRO}, enquanto o
comportamento do motor
(ativações, desativações e conflitos) é idêntico nos dois hospedeiros. Com
telemetria física e erro de posição, o \texttt{cooldown} reduziu as transições
espúrias em 11\,\% a 22\,\%, sem eliminá-las. O objetivo geral foi atingido
para o cenário avaliado, nos dois hospedeiros.

Entre os desafios encontrados, destacam-se: manter a gramática LL(1) admitindo
tanto referências simples (\texttt{frame\_budget}) quanto qualificadas
(\texttt{view "FPV".scale}); incorporar a unidade ao \textit{token} numérico
sem ambiguidade com identificadores (\texttt{10 frames} contra
\texttt{10 mode}); lidar com as consequências da convenção léxica de
enumerados em caixa alta, que obriga nomes de regra em PascalCase e exige
validar o modelo do hospedeiro contra as palavras reservadas; decidir V6 e V7
de forma exata sem recorrer a um resolvedor externo; definir a semântica
temporal do \texttt{ASSERT}, verificado no quadro seguinte ao da aplicação do
plano; separar a medição de custo da própria renderização, o que levou ao
relógio por carga calibrado; e tornar reproduzíveis experimentos que dependem
de um simulador físico, o que levou a avaliar sobre gravações versionadas.

Como expansão, propõem-se: medir o custo sobre um compositor em GPU;
exercitar os pontos de junção em execução e desenhar \texttt{HIGHLIGHT} e
\texttt{MARK}; acrescentar termo fixo por passo ao relógio por carga e
calibrar o \texttt{SuitabilityPass} nas demais superfícies; fechar o laço entre
a política e o voo no CoppeliaSim; avaliar períodos de \texttt{cooldown}
maiores e filtragem da percepção contra o erro de posição; declarar no modelo a
telemetria dos outros
cinco cenários excepcionais; estender V4 a reduções de resolução e de
frequência de superfícies críticas; verificar exaustivamente os até $2^n$
planos alcançáveis de um programa; oferecer uma janela de visualização das
gravações, como ferramenta de demonstração; aceitar identificadores Unicode e
números negativos; e conduzir o estudo de legibilidade com agrônomos e
pilotos~\cite{rodrigues2017}.

%=======================================================================
\begin{thebibliography}{00}

\bibitem{mernik2005} M. Mernik, J. Heering e A. M. Sloane, ``When and how to
develop domain-specific languages,'' \textit{ACM Computing Surveys}, vol. 37,
n. 4, pp. 316--344, 2005.

\bibitem{vandeursen2000} A. van Deursen, P. Klint e J. Visser, ``Domain-specific
languages: an annotated bibliography,'' \textit{ACM SIGPLAN Notices}, vol. 35,
n. 6, pp. 26--36, 2000.

\bibitem{nascimento2012} L. M. do Nascimento, D. L. Viana, P. A. M. S. Neto,
D. A. O. Martins, V. C. Garcia e S. R. L. Meira, ``A systematic mapping study on
domain-specific languages,'' in \textit{Proc. 7th Int. Conf. on Software
Engineering Advances (ICSEA)}, 2012, pp. 179--187.

\bibitem{rodrigues2017} I. Rodrigues, M. de Borba Campos e A. F. Zorzo,
``Usability evaluation of domain-specific languages: a systematic literature
review,'' in \textit{Human-Computer Interaction (HCI International 2017)}, LNCS
vol. 10271, Cham: Springer, 2017, pp. 522--534.

\bibitem{boukham2025} H. Boukham, K. Younsi Dahbi e D. Chiadmi,
``Domain-specific languages for algorithmic graph processing: a systematic
literature review,'' \textit{Algorithms}, vol. 18, n. 7, art. 445, 2025.

\bibitem{breuker2014} D. Breuker, ``Towards model-driven engineering for big data
analytics -- an exploratory analysis of domain-specific languages for machine
learning,'' in \textit{Proc. 47th Hawaii Int. Conf. on System Sciences (HICSS)},
Waikoloa, HI, EUA, 2014.

\bibitem{alves2022} L. Alves, J. D. Pereira, N. Aragão, M. Chagas e P. H. M.
Maia, ``DRESS-ML: a domain-specific language for modelling exceptional scenarios
and self-adaptive behaviours for drone-based applications,'' in \textit{Proc.
44th Int. Conf. on Software Engineering: Software Engineering in Society (ICSE
SEIS)}, Pittsburgh, PA, EUA, 2022, pp. 56--66.

\bibitem{cheng2012} S.-W. Cheng, D. Garlan e B. Schmerl, ``Stitch: a language for
architecture-based self-adaptation,'' \textit{Journal of Systems and Software},
vol. 85, n. 12, pp. 2860--2875, dez. 2012, doi: 10.1016/j.jss.2012.02.060.

\bibitem{aho2007} A. V. Aho, M. S. Lam, R. Sethi e J. D. Ullman,
\textit{Compilers: Principles, Techniques, and Tools}, 2nd ed. Boston, MA, EUA:
Pearson/Addison-Wesley, 2007.

\bibitem{hopcroft2007} J. E. Hopcroft, R. Motwani e J. D. Ullman,
\textit{Introduction to Automata Theory, Languages, and Computation}, 3rd ed.
Boston, MA, EUA: Pearson/Addison-Wesley, 2007.

\end{thebibliography}

\end{document}